% !TeX program = xelatex
\documentclass[12pt,a4paper]{ctexart}
\usepackage{geometry}
\usepackage{fontspec}
\usepackage{fancyhdr}
\usepackage{graphicx}
\usepackage{amssymb}
\usepackage{amsmath}
\usepackage{booktabs}
\usepackage{multirow}
\usepackage{array}
\usepackage{tabularx}
\usepackage{caption}
\usepackage{placeins}
\usepackage{float}
\usepackage{listings}
\usepackage[table]{xcolor}
\usepackage{xeCJK}
\geometry{left=2.5cm,right=2.5cm,top=2.5cm,bottom=2.5cm,footskip=16pt}
\linespread{1.3}
\graphicspath{{figures/}}
\pagestyle{fancy}\fancyhf{}\renewcommand{\headrulewidth}{0pt}\cfoot{\zihao{-5}\thepage}
\setmainfont{Times New Roman}
\setCJKmainfont{SimSun}[AutoFakeBold=2.5]
\setCJKfamilyfont{zhhei}{SimHei}[AutoFakeBold=2.5]
\renewcommand{\heiti}{\CJKfamily{zhhei}}
% ①②③ 等带圈数字在 TNR 无字形，声明其走中文字体
\xeCJKDeclareCharClass{CJK}{"2460 -> "24FF, "00A0 -> "00BF, "0370 -> "03FF, "2010 -> "2027, "2070 -> "209F, "2190 -> "21FF, "2200 -> "22FF, "27F0 -> "27FF, "2A00 -> "2AFF}
\setlength{\parindent}{2\ccwd}
\ctexset{
  section={name={,、},number={\chinese{section}},format={\zihao{4}\heiti\centering},},
  subsection={format={\zihao{-4}\heiti\raggedright}},
  subsubsection={format={\zihao{-4}\heiti\raggedright}}
}
\lstdefinestyle{codeblock}{
  basicstyle=\ttfamily\zihao{5}, breaklines=true, columns=fullflexible,
  keepspaces=true, showstringspaces=false, frame=single, numbers=none,
  xleftmargin=1em, xrightmargin=0.11cm, aboveskip=6pt, belowskip=6pt,
  extendedchars=true, inputencoding=utf8,
  literate={∈}{{$\in$}}1 {≤}{{$\le$}}1 {≥}{{$\ge$}}1 {←}{{$\leftarrow$}}1 {∅}{{$\varnothing$}}1 {θ}{{$\theta$}}1 {ρ}{{$\rho$}}1 {∩}{{$\cap$}}1 {°}{{$^\circ$}}1 {±}{{$\pm$}}1 {·}{{$\cdot$}}1 {¹}{{$^{1}$}}1 {×}{{$\times$}}1 {Δ}{{$\Delta$}}1 {Σ}{{$\Sigma$}}1 {β}{{$\beta$}}1 {π}{{$\pi$}}1 {⁻}{{$^{-}$}}1 {→}{{$\rightarrow$}}1 {⇒}{{$\Rightarrow$}}1 {⇔}{{$\Leftrightarrow$}}1 {⟺}{{$\Longleftrightarrow$}}1 {-}{{$-$}}1 {∪}{{$\cup$}}1 {≈}{{$\approx$}}1
}
\captionsetup{font={small,bf},labelsep=quad}
\newcommand{\upcite}[1]{\textsuperscript{\cite{#1}}}
% 附录代码清单样式：小五号、单倍行距、带边框
\lstdefinestyle{codeappx}{
  basicstyle=\ttfamily\zihao{6}, breaklines=true, columns=fullflexible,
  keepspaces=true, showstringspaces=false, frame=single, numbers=none,
  xleftmargin=1em, xrightmargin=0.11cm, aboveskip=4pt, belowskip=4pt,
  extendedchars=true, inputencoding=utf8,
  literate={⁻}{{$^{-}$}}1 {¹}{{$^{1}$}}1
}
\begin{document}

\begin{center}
	\textbf{\zihao{3}\heiti 基于覆盖证书与顺路清除的干扰源全域保证检出与快速清除}
\end{center}
\vspace{-6pt}
\heiti\zihao{4}\begin{center}摘要\end{center}
\vspace{-6pt}
\songti\zihao{-4}\par
本文研究机器狗携带测向机在半径 1800 m 圆形区域内自动搜索、定位并清除干扰源的问题。全文以“一条示向度 $\rightarrow$ 后验可行集 $\rightarrow$ 主动选点 $\rightarrow$ 后验更新 $\rightarrow$ 清除判据 $\rightarrow$ 开火或补测”为主线，四问共用同一定位内核（示向带取交、多示向度一致性筛选、分量中位数初值、{\heiti Huber 加权最小二乘}）与同一清除判据 $\rho\le20$ m：先用可行性证书回答“是否可能漏源”，再用限时巡回与顺路清除回答“如何尽快处置”。

针对问题一，建立楔形交会—凸集直径—{\heiti 最小外接圆}模型：以半平面裁剪精确构造定位区域，用旋转卡壳求直径、以 Welzl 算法求{\heiti 最小外接圆}，并以 Thales 判据给出“以直径 $D$ 为直径的圆能否覆盖该区域”的等价判定。结果表明该圆并不具有普遍覆盖保证：2000 例两点交会构型中不满足覆盖者占 1.50\%（Wilson 95\% 置信区间 [1.05\%, 2.13\%]），反例族可将比值推至 Jung 上界 $2/\sqrt3\approx1.1547$，而求解器与暴力枚举逐位一致（偏差 $1.137\times10^{-13}$ m）。

针对问题二，将其建模为带接收约束的{\heiti 主动观测点设计} $\max\,\mathbb{E}[\ln(\rho_1/\rho_2)]$：状态取一次观测确定的后验可行集，约束为对后验集中所有可能源位置均保证可收信号。最优第二点位于示向线上 840 m、侧向偏移 $-$543 m 处，交会角 56.4°，期望几何收缩效用 3.302 nat，二次测量后可直接开火的概率由贴近源方案的 4.2\% 提高到 28.3\%；退到保证可收的可执行解，定位区域直径仍缩小 22.8 倍。

针对问题三，由“盲区最大值等于环半径”与“巡回里程 $\approx2\pi R$”两条结构性结论把布局定在约束边界，取单环 990 m$\times$15 点建证集，用“频道 $\times$ 空间”{\heiti 覆盖证书}替代启发式终止，并以绕道代价判据实现{\heiti 顺路清除}，巡回里程 6753 m。在最终配置下 16 局演练全部 100\% 清除（平均 364.6 s/源）之后，完成三次正式测试：三局均 100\% 清除（合计 35/35），平均定位清除时间 395.0 s/源（逐局 331.9～435.7 s/源），三局均由覆盖证书终止而非航点耗尽。

针对问题四，定向源的负观测不含距离信息，终止依据重建为“位置 $\times$ 朝向”空间的全域保证检出证书（31 点，另配 50 点补给），在 5 m$\times1^\circ$ 网格 1.465 亿个位置-朝向对上零漏覆盖；“去掉外环”的反证实跑漏检 4 源，可见外环结构不可省。最终配置下 18 局演练：16 局 100\% 清除（另 2 局为反证与压力实验），平均定位清除时间主体为约 400～670 s/源。三次正式测试全部 100\% 清除（合计 37/37），平均定位清除时间 568.5 s/源（逐局 504.1～620.3 s/源），共 49 次清除调用中 37 次命中，即每个源恰好清除一次；第 2 局的两个弱观测源分别由主动补测与收尾清除兜底完成，验证了长尾修复的有效性。

本文对覆盖口径、接收半径、样本量与参数扰动作了稳健性检验，四问结论均未改变；“{\heiti 覆盖证书} + {\heiti 顺路清除}”框架可迁移到频谱监测、定向天线对准、航点巡检等“带朝向约束的可达性 + 巡回清除”类问题。

\vspace{6pt}
\noindent{\heiti\textbf{关键词：}}干扰源定位与清除；覆盖证书；顺路清除；航点布局优化；最小外接圆；主动观测点设计；Huber 加权最小二乘
\songti
\newpage

\FloatBarrier

\section{问题重述}

\subsection{问题背景}

无线电干扰源的人工排查需技术人员携带测向机徒步进入目标区域，逐点观察场强与方向并凭经验逼近，效率不高且在危险环境中难以持续\upcite{ref:Li2025}。本题要求改用搭载测向机的机器狗，在半径 1800 米的圆形区域内自动完成搜索、定位与清除。测向机测得的是示向度\upcite{ref:Torrieri}，只给方位、不给距离；误差在同一地点固定，故重复检测无法改善。干扰源分全向与定向两类，后者仅在朝向两侧各 90 度的扇区内辐射，这一方位盲区使“测不到”并不等于“不存在”。题目要求确保全部清除与耗时尽可能短，二者相互制约：排查越疏，漏检风险越大；排查越密，耗时越长。

\subsection{问题要求}

问题一：已知若干检测点的坐标及某干扰源关于这些检测点的示向度，要求给出用交会定位法计算定位区域直径的算法（定位区域直径指该区域内任意两点距离的最大值），并判断以该直径为直径的圆能否覆盖该区域，给出判定依据。

问题二：已知某全向干扰源在一个检测点处测得的示向度，要求确定第二个检测点的选择策略，使定位效果较好，并给出候选区域。

问题三：目标区域内均为全向源，总数在 10 至 16 之间但具体个数未知；机器狗从原点出发、测向机初始频道为 1，要求制定自动搜索、定位与清除的策略及算法，在确保全部清除的前提下使耗时尽可能短。检验方式为模拟器演练，给出被清除源比例与平均定位清除时间，并在演练后完成三次正式测试，按题面表 1 报告。

问题四：条件与问题三相同，但区域内同时含全向与定向源，定向源的个数与朝向未知，要求给出检测与清除的策略及算法，检验与报告方式同问题三。

\FloatBarrier

\section{问题分析}

\subsection{问题一的分析}

问题一本质是凸集构造与最小外接圆计算：示向度连同误差界给出楔形，各楔形取交、再与目标圆域求交即得定位区域；该区域为凸集，直径必由一对顶点取得，可用旋转卡壳精确求解。判定覆盖则引入最小外接圆——以直径 $D$ 为直径的圆覆盖该区域当且仅当 $\rho_{\min}=D/2$；由 Jung\upcite{ref:Jung} 定理 $\rho_{\min}/(D/2)\in[1,\ 2/\sqrt3]$，故该圆对一般区域{\heiti 不具有普遍覆盖保证}。本问的落点在清除判据：最小外接圆半径才是清除可行性的判据（$\rho_{\min}\le20$ m 时其圆心即保证清除的开火点），问题一由此接入问题三、四。

\subsection{问题二的分析}

问题二不是静态几何求最优点，而是带接收约束的主动观测点设计：第二点的价值取决于它能把第一次观测给出的后验可行集压缩多少，且须保证观测能够完成。困难根源在于“有效接收半径介于 1000 与 1500 米”是源自身的接收能力，由检出只能推知距离的{\heiti 上界}，故单条示向度不含径向信息，后验集是一条很长的楔形。选点目标取后验集的期望收缩（以最小外接圆半径为不确定性度量），约束是“保证收到”——对后验集中每个可能源位置都满足到第二点距离不超过 1000 米；该约束集一般不是圆盘，用圆盘作外近似对保证型优化不安全。

\subsection{问题三的分析}

问题三本质是覆盖约束下的序贯决策问题，难点在三层。几何层须区分“至少一个点测到”与“能交会定位”：前者仅需少数航点即可覆盖全域，但源在圆心附近时只有 1 个航点可测、无法交会，故须加严为二重可见，并由此引出“频道 $\times$ 空间”覆盖证书。估计层的困难在于：单条示向线只含方位不含径向，两线交会误差被交会角正弦的倒数放大\upcite{ref:Chen2025}，同点偏差固定又使重复检测无增益。决策层则在覆盖有保证的前提下安排动作次序，以虚拟时间最小为目标，“连续若干步无新发现”仅为非保证型的加速规则。

\subsection{问题四的分析}

问题四在问题三基础上增加观测可达性约束：定向源仅在朝向两侧各 90 度的扇区内辐射，故某点测不到它可能仅因该点落在扇区之外；对定向源的负观测不含距离信息，其排除逻辑不能与全向源混用。本文因而利用题面允许提交域外坐标的条件在区域外增设环绕带，并把覆盖判据分三级逐层加严：一重可见（任意朝向下至少一个航点可见，等价于可测航点集方位排序后最大相邻间隙不超过 180 度）、二重可见（至少两个航点可见，为交会定位前提）、几何可定位（存在可见点对且交会角不病态）。三者均为几何充分条件，故统称“证书”；其全域验证口径与所留余量在第八章交代。

上述四问在方法上构成一条闭环，其整体脉络可描述为：由一次示向度出发，依次串起后验可行集、主动选点、后验更新、清除判据与开火或补测的决策。


\FloatBarrier

\section{模型假设}

题面已给定的条件（频道编号、源类型划分与覆盖范围、接收半径区间、示向度误差有界且同点固定、平面简化、计时与清除判定等）本文直接采用。补充假设如下。

{\heiti 假设 1（单源单频道）}　每个频道至多对应一个未清除源，不同源信号互不影响，故多个测点示向度求得的交会点唯一对应该频道的源。

{\heiti 假设 2（源静止）}　单次测试过程中，各源的位置与朝向不随时间变化。

{\heiti 假设 3（保守接收半径）}　凡涉及可探测性的保证性结论，一律取接收半径下界 1000 米为可接收距离，使结论对任意源成立；该取值是保证型设计口径，并非对真实半径的假定。

{\heiti 假设 4（坐标提交精确）}　机器狗可准确到达提交坐标，接口返回的位置即真实位置，从而忽略转向、加减速与到达时刻的偏差。

{\heiti 假设 5（不引入误差分布假设）}　同一测点的误差固定、不同测点的误差互不相关，故更换测点可获取新信息；题面只给出误差的界而无分布，故本文不以正态假设合并多测线——保证型结论一律以“角度区间取交”给出，仅点估计与排序使用稳健加权最小二乘。

{\heiti 假设 6（模拟器可信）}　官方模拟器按题面规则运行，其计时与检测、清除反馈真实反映当前状态。



\FloatBarrier

\section{符号说明}

全文主要符号如表~\ref{tab:symbols} 所示，同一符号含义全文一致。角度按题面约定，取正东方向为 0 度，逆时针为正，范围为 0 度至 360 度（含 0 度、不含 360 度）；方位角与示向度均沿用此约定。

\begin{table}[htbp]
\centering
\caption{符号说明}\label{tab:symbols}
{\zihao{6}\renewcommand{\arraystretch}{0.92}\begin{tabular}{@{}c >{\raggedright\arraybackslash}p{5.05cm} c @{\hspace{0.7em}} c >{\raggedright\arraybackslash}p{5.05cm} c@{}}
\toprule
符号 & 含义 & 单位 & 符号 & 含义 & 单位 \\
\midrule
$O$ & 目标区域圆心，取为坐标原点 & m & $v$ & 机器狗移动速度，取 5 & m/s \\
$R$ & 目标区域半径，取 1800 & m & $n_{\rm src}$ & 干扰源总数，介于 10 与 16 之间 & 个 \\
$S_k$ & 第$k$ 个检测点的位置 & m & $n_{\rm cleared}$ & 已被清除的干扰源个数 & 个 \\
$\theta_k$ & 检测点$S_k$ 处的示向度 & 度 & $n_{\rm meas},\,n_{\rm sw}$ & 检测次数、相邻两次检测频道不同时的切频次数 & 次 \\
$\varepsilon$ & 示向度误差界，取 1 & 度 & $n_{\rm ok},\,n_{\rm miss}$ & 清除成功、失败的次数 & 次 \\
$W_k$ & 由示向度与误差界确定的楔形区域 & — & $T_{\rm virt}$ & 虚拟时间，上限 360000 & s \\
$\mathcal P_G$ & 精确示向度下的定位区域（楔形交会区域与目标圆域之交） & — & $T_{\rm real}$ & 现实时间，即程序运行时间，上限 1200 & s \\
$\mathcal T$ & 含系统偏差时的实际可行源集 & — & $T_{\rm clear}$ & 定位清除总时间（按虚拟时间计） & s \\
$D$ & 定位区域的直径 & m & $\bar T_{\rm clear}$ & 平均定位清除时间，$T_{\rm clear}$ 除以被清除源个数 & s \\
$\rho_{\min}$ & 集合的最小外接圆半径 & m & $\Omega_1$ & 第一次观测后的后验可行集（楔形 与 距离上界 与 目标圆域） & — \\
$\psi$ & 两条方位线在源处的交会角 & 度 & $\ell$ & 第二检测点沿第一条示向线的前进量 & m \\
$S_2$ & 第二次观测所选检测点的位置 & m & $g$ & 第二检测点相对第一条示向线的侧向偏移 & m \\
$\mathcal S_2$ & 第二次观测点的候选区域 & — & $a$ & 第二检测点到源的垂距 & m \\
$\mathcal T_2$ & 第二次观测后收缩得到的可行集 & — & $R_{\rm geo}$ & 可达域 & — \\
$\Delta_{\max}$ & 观测点方位排序后的最坏角隙 & 度 & $R_{\rm recv}$ & 保证可收集，即对后验集中所有可能源位置成立的接收可行域 & — \\
$r_i$ & 第$i$ 个干扰源的有效接收半径 & m & $R_{\rm good}$ & 交会角非病态条件集 & — \\
$r_{\min}$ & 有效接收半径下界，取 1000 & m & $P_x$ & 源位于$x$ 时其下界半径邻域内的航点集合 & — \\
$d_{\rm clear}$ & 清除半径，取 20 & m & $P_c$ & 某频道上已检测且无信号的航点集合 & — \\
$d_1,\,d_2$ & 基准场景中第一、二次观测点到源的距离 & m & $K^{*}$ & 全域保证检出证书集的最少点数 & 个 \\
$G$ & 基准场景中的源位置 & m & & & \\
\bottomrule
\end{tabular}}
\end{table}

\FloatBarrier


\section{问题一模型的建立与求解}

\subsection{思路流程}

问题一要从检测点坐标与示向度出发，算出定位区域的直径，并判断以该直径为直径的圆能否覆盖该区域。本章的求解分四步：先由各检测点连同正负 1 度误差界确定楔形，各楔形取交、再与目标圆域求交，得到定位区域；再求该区域的直径与最小外接圆；然后按 Thales 判据判定覆盖是否成立；最后把最小外接圆半径接入清除判据，作为能否一次清除的依据。其中半平面裁剪是精确的，圆盘边界以足够细的多边形近似并给出径向误差上界。


\subsection{数据预处理与可视化}

本章输入为题面给定的检测点坐标与示向度。为检验算法自洽性与覆盖判据表现，另造三组数据：规则构型用于一致性校验与构型扫描；蒙特卡洛\upcite{ref:MC}采样（源位在半径 1700 米内均匀，两测点到源距离取 400～1500 米与 250～800 米，交会角 20～160 度）共 2000 个两点交会构型，用于统计直径与覆盖比例；反例族（三点按 120 度均布于半径 1000 米圆周、示向度推向楔形边界）用于构造最坏构型。口径上，坐标按 9 位小数去重，判据比较取 $10^{-9}$ 容差，并附一份无容差、含浮点平局的结果作对照。

可视化见图~\ref{fig:q1-geometry-cover}（左：定位区域与最小外接圆；右：可覆盖与不可覆盖对照）、图~\ref{fig:q1-counterexample}（反例族逼近 Jung 上界）；判据口径敏感性与尾部互补分布见附录图 D3。

\begin{figure}[htbp]
\centering
\begin{minipage}[t]{0.46\linewidth}\centering
\includegraphics[width=\linewidth]{fig_q1_geometry.png}
\end{minipage}\hfill
\begin{minipage}[t]{0.46\linewidth}\centering
\includegraphics[width=\linewidth]{fig_q1_cover_compare.png}
\end{minipage}
\caption{左：交会定位区域、直径与最小外接圆；右：以直径为直径的圆能否覆盖定位区域}\label{fig:q1-geometry-cover}
\end{figure}


\begin{figure}[htbp]
\centering
\includegraphics[width=0.42\linewidth]{fig_q1_counterexample.png}
\caption{反例族随示向度偏移逼近 Jung 上界}\label{fig:q1-counterexample}
\end{figure}

\FloatBarrier

\subsection{模型建立}

（1）定位区域的构造。设第 $k$ 个检测点为 $S_k$、示向度为 $\theta_k$、误差界为 $\varepsilon$，定义半角 $\varepsilon$ 的闭凸楔形

\begin{equation}
	\text{楔形} W_k=\{S_k+\rho\,u(\varphi):\rho>0,\ |\varphi-\theta_k|\le\varepsilon\}
	\label{eq:wedge}
\end{equation}

则示向度交会区域为各楔形之交，与题面图 2 红色四边形一致。题面规定所有干扰源位于半径 1800 米圆域内，故本文定位区域取

\begin{equation}
	\text{定位区域} \mathcal P_G=(\bigcap_k W_k)\cap\bar B(O,1800)
	\label{eq:region}
\end{equation}

加入目标圆域有两条理由：一是圆域是题面已给的信息；二是当检测点与方向组合近共线时，纯示向交会区域可能无界、直径与覆盖判据失去意义，与圆域求交后区域恒有界，问题才适定。

（2）直径与最小外接圆。记

\begin{equation}
	\text{直径} D=\max_{X,Y\in\mathcal P_G}\|X-Y\|
	\label{eq:diameter}
\end{equation}

该区域是有限个半平面与圆盘之交，为凸紧集，故直径与最小外接圆半径都存在（由顶点或圆弧端点确定）。

（3）覆盖判据。以 $D$ 为直径的圆覆盖该区域当且仅当

\begin{equation}
	\rho_{\min}=D/2\iff\exists(A,B),\ \|A-B\|=D,\ (X-A)\cdot(X-B)\le0,\ \forall X\in\mathcal P_G
	\label{eq:criterion}
\end{equation}

后一式即 Thales 判据，即存在一对实现直径的点 $A$、$B$，使区域整体落进以 $AB$ 为直径的闭圆盘。

（4）理论界。由 Jung 定理

\begin{equation}
	D/2\le\rho_{\min}\le D/\sqrt3\text{，即} \rho_{\min}/(D/2)\in[1,\ 2/\sqrt3=1.1547]
	\label{eq:jung}
\end{equation}

下限对应“某个直径圆恰好覆盖”，上限意味着任何半径小于该值的圆盖不住；最坏情形由近正三角形三点达到。

（5）清除判据。题面附录 2(8) 规定精确定位与清除仅取决于清除点到源的距离，阈值为 20 米，故保证成功的开火点当且仅当

\begin{equation}
	\rho_{\min}(\mathcal P_G)\le d_{\rm clear}=20 \text{米}
	\label{eq:clear-threshold}
\end{equation}

这时最小外接圆圆心即该开火点。

\subsection{模型求解}

求解算法如下：

\begin{lstlisting}[style=codeblock]
输入：观测集 {(S_k, theta_k)}，eps=1°，R=1800
1  P <- 圆盘(O,R) 的多边形逼近                 # N=720（R=1500 m），径向偏差 0.0143 m
2  for each (S,theta):
       P <- clip_halfplane(P, theta-eps 边界)        # 半平面裁剪，精确
       P <- clip_halfplane(P, theta+eps 边界)
3  V <- convex_hull(P)                          # Andrew 单调链
4  D <- rotating_calipers(V)                    # 旋转卡壳，O(m)
5  (C,rho_min) <- Welzl(V)                        # 最小外接圆，期望线性
6  if |rho_min - D/2| <= 1e-9: 返回“以 D 为直径的圆可覆盖”
   else: 返回“不能覆盖”，并输出漏出直径圆的顶点
\end{lstlisting}

其中第 2 步半平面裁剪是精确的；与圆盘求交由第 1 步多边形逼近承担，径向误差 $R(1-\cos\frac{\pi}{N})$ 在 $R=1500$ 米、$N=720$ 时为 0.0143 米。

求解结果如下（详见附录表 D10）：（a）算法一致性——旋转卡壳与暴力枚举在 400 组随机构型上逐位一致，偏差在浮点精度极限内；（b）反例族——三点按 120 度均布，当示向度推至楔形边界极限时，比值达到 Jung 上界 $2/\sqrt3$（1.154701），说明该上界可达；（c）蒙特卡洛统计——2000 个两点交会构型下，容差判据不满足覆盖的比例为 1.50\%，其 Wilson 95\%\upcite{ref:Wilson} 置信区间由式~\eqref{eq:wilson} 给出，为 [1.05\%, 2.13\%]，改用无容差的浮点平局口径升至 11.35\%（仅作口径对照），独立重抽 2000 例复现 1.65\%；（d）极值搜索——2 点 20 万例、3～6 点各 6 万例的比值上界见附录表 D10；（e）口径敏感性——主模型 $\mathcal P_G$ 与“纯示向楔形区域”在本采样下数值完全一致，说明该 2000 例中目标圆域未发生过截断（仅为样本结论）。

\begin{equation}\label{eq:wilson}
  \frac{\hat p+\dfrac{z^2}{2n}\pm z\sqrt{\dfrac{\hat p(1-\hat p)}{n}+\dfrac{z^2}{4n^2}}}{1+\dfrac{z^2}{n}},\qquad \hat p=\frac{k}{n},\ z=1.96
\end{equation}



\subsection{结果分析与讨论}

本文结论是，以定位区域直径为直径的圆不具有普遍覆盖保证——只有当最小外接圆恰好以一对径点为直径端点时（式~\eqref{eq:criterion}）才覆盖，否则至少漏掉一个顶点。该结论由反例驱动，而非由高频现象驱动：本文采样下不满足覆盖的构型仅占 1.50\%，但确实存在能把比值推至 Jung 上界的反例族，故本文同时给出解析反例与统计比例并附置信区间，并说明该比例依赖于采样分布。

判据口径必须写明：若不设容差，$10^{-9}$ 量级浮点平局会使比例由 1.50\% 升至 11.35\%，本文以容差判据为结论，以无容差结果作口径对照。


另有两处衔接与局限：本章落点是清除判据——直径只是尺度指标，最小外接圆半径才是“能否一次清除”的判据（式~\eqref{eq:clear-threshold}），问题一结果由此成为问题三、四决策层的输入；又因本章结论建立在“检测点与示向度精确已知”上，在线场景须按问题三的做法把定位区域视为可行源集并逐点更新（详见第七章）。

\FloatBarrier

\section{问题二模型的建立与求解}

\subsection{思路流程}

问题二要在只掌握一条示向度的条件下确定第二次测量点的位置，使定位效果尽可能好。本章把它建模为带接收约束的主动观测点设计：状态取第一次观测确定的后验可行集而非单点估计，目标为第二次观测对不确定性度量的期望收缩，约束为必须保证能收到信号。求解分三步：按前进量与侧向偏移枚举候选点，逐个计算交会区域的最小外接圆半径以评估目标；再以不依赖先验的最坏交会角准则作对照；最后给出在线执行方式及其与清除判据的衔接。


\subsection{数据预处理与可视化}

本章输入由构造给出：后验可行集由一条示向度、误差界、距离上界 1500 米与目标圆域求交得到，作为采样先验的支撑集；候选点网格在示向线局部坐标系下按前进量与侧向偏移划分，先用解析误差尺度粗筛、再对排序靠前者作精确几何计算；先验样本在可行集内按面积均匀抽取 120 个源位置。

两处口径须说明。其一，目标函数含对数项，对半径取下限保护；接近零的样本须保留并按该下限处理，不宜剔除，否则引入选择偏差。其二，第二测点落在源 5 米内时，按题面附录 2(9) 返回近场结果，不含示向度，应视为“可直接清除”并计入成功。

可视化见图~\ref{fig:q2-geometry-scale}（左：第二检测点的最优几何与目标函数平坦性；右：解析误差尺度场与精确直径场的一致性）、图~\ref{fig:q2-candidate}（保证可收集与候选区域三层）、图~\ref{fig:q2-flatness-infogain}（左：$D(a)$ 的双段律；右：后验几何收缩效用场与三类选点对比）。

\begin{figure}[htbp]
\centering
\begin{minipage}[t]{0.42\linewidth}\centering
\includegraphics[width=\linewidth]{fig_q2_geometry.png}
\end{minipage}\hfill
\begin{minipage}[t]{0.42\linewidth}\centering
\includegraphics[width=\linewidth]{fig_q2_scale_vs_exact.png}
\end{minipage}
\caption{左：第二检测点的最优几何与目标函数平坦性；右：解析误差尺度场与精确直径场的一致性}\label{fig:q2-geometry-scale}
\end{figure}

\begin{figure}[htbp]
\centering
\includegraphics[width=0.46\linewidth,height=0.27\textheight,keepaspectratio]{fig_q2_candidate_region.png}
\caption{保证可收集与候选区域}\label{fig:q2-candidate}
\end{figure}

\begin{figure}[htbp]
\centering
\begin{minipage}[t]{0.42\linewidth}\centering
\includegraphics[width=\linewidth]{fig_q2_flatness.png}
\end{minipage}\hfill
\begin{minipage}[t]{0.42\linewidth}\centering
\includegraphics[width=\linewidth]{fig_q2_infogain.png}
\end{minipage}
\caption{左：$D(a)$ 的双段律（拐点 $a^{*}=d_1\tan1^\circ$ 与平台段）；右：后验几何收缩效用场与三类选点对比}\label{fig:q2-flatness-infogain}
\end{figure}

\FloatBarrier

\subsection{模型建立}

（1）后验可行集。一次示向度只给出方位，结合距离上界与目标圆域得

\begin{equation}
	\Omega_1=(W_1\cap\bar B(S_1,1500))\cap\bar B(O,1800)
	\label{eq:omega1}
\end{equation}

由一次成功检出只能推知距离不超过 1500 米，故 $\Omega_1$ 是一条径向长度可达 1500 米的狭长楔形，单条示向度不含径向信息。

（2）选点准则。选定第二点 $S_2$ 后，第二次示向度把 $\Omega_1$ 收缩为 $\mathcal T_2$。以最小外接圆半径为不确定性度量，准则为

\begin{equation}
	\max_{S_2}\ \mathbb{E}_{x\sim p(x\mid\theta_1)}\big[\log\big(\rho_{\min}(\Omega_1)/\rho_{\min}(\mathcal T_2(x;S_2))\big)\big]
	\label{eq:utility}
\end{equation}

该式仅对源位置取期望且将第二次观测视为无噪声，故记为“后验几何收缩效用”而非期望信息增益；题面未给先验，取 $\Omega_1$ 上均匀分布为最大熵基线。

（3）接收可行性约束。要保证第二次观测一定收到信号，须对 $\Omega_1$ 中所有可能源位置同时满足到第二点距离不超过 1000 米，即

\begin{equation}
	R_{\rm recv}=\bigcap_{x\in\Omega_1}\bar B(x,r_{\min})=\bigcap_{v\in\operatorname{ext}(\Omega_1)}\bar B(v,r_{\min})
	\label{eq:recv}
\end{equation}

第二个等号成立的依据是：距离关于源位置为凸函数，其最大值必在凸集的极点取得；$\Omega_1$ 的极点只有少数几个，故 $R_{\rm recv}$ 是至多六个圆盘的交，一般不是圆盘。候选区域为

\begin{equation}
	\mathcal S_2=R_{\rm geo}\cap R_{\rm recv}\cap R_{\rm good}
	\label{eq:candidate}
\end{equation}

其中 $R_{\rm geo}$ 为可达域、$R_{\rm good}$ 为交会角非病态条件集。

（4）几何参照：直径的双段律。若把第二点取在过源的垂线上、垂距为 $a$，则两点交会区域的直径满足

\begin{equation}
	D(a)=d_1\tan1^\circ+\min(a,\ d_1\tan1^\circ)
	\label{eq:twoseg}
\end{equation}

其拐点为 $a^{*}=d_1\tan1^\circ$，超过拐点后进入与 $a$ 无关的平台段。由此可知，“越靠近源越好”仅当能足够精确地掌握源位置时才成立。

\subsection{模型求解}

求解算法如下：

\begin{lstlisting}[style=codeblock]
输入：S1, theta1, eps=1°, r_min=1000, r_max=1500
1  Omega1 <- 楔形  ∩  圆盘(S1,r_max)  ∩  圆盘(O,1800)          # 多边形裁剪
2  R_recv <- 对 Omega1 各极点 v 的圆盘(v, r_min) 取交               # 至多 6 个圆盘之交
3  for each 候选 S2  in  R_geo  ∩  R_recv  ∩  R_good:
       采样 x ～ U(Omega1)
       对每个 x: theta2 <- bearing(S2,x); T2 <- Omega1  ∩  楔形(S2,theta2)
                 rho2 <- Welzl(T2)
       U(S2) <- mean( log(rho1 / max(rho2,eps)) )
4  S2* <- argmax U(S2)；并另算最坏交会角准则的最优解作对照
\end{lstlisting}

求解结果（基准场景 $S_1=(0,0)$、$G=(0,1200)$、$d_1=1200$ 米、$\theta_1=90$ 度）：

（a）单次观测的定位区域直径已接近接收半径上界（数值见附录表 D7）。

（b）期望效用最优解与三类方案对照见表~\ref{tab:q2-utility}。效用最优解位于示向线上 840 米、侧向偏移 $-$543 米，交会角约 56 度，其二次观测后可直接开火的概率（28.3\%）远高于其余三种方案（见表~\ref{tab:q2-utility}）。

（c）基准场景的几何量详见附录表 D7。双段律的拐点、平台值及各垂距对应的直径等逐项数值均见附录表 D7，其中源位已知时的最优直径相对单次缩小 47.7 倍，可作几何上界参照。

（d）多场景鲁棒性（24 个随机场景，$d_1$、交会角与源位均随机）：最优垂距中位数 31 米（范围 5～89 米），交会角中位数 89.6 度（范围 72.4～101.9 度），逐场景明细见支撑材料。

（e）解析式与精确几何的关系：解析误差尺度给出的最优点位置与精确几何完全一致（偏差 0.0 米 / 0.0 度），但其幅度仅具缩放律性质（全网格相关系数约 0.82）。故正文口径为：解析式用于候选点排序，精确多边形几何给出最终直径。

（f）候选区域取 $R_{\rm recv}$ 与可达域、交会角非病态条件集的交集；其中 $R_{\rm recv}$ 以圆盘作外近似时存在约 26.2 米的横向偏差上界，对保证型优化不安全（见 6.3（3））。

\begin{table}[htbp]
\centering
\caption{四类选点方案的后验几何收缩效用对比}\label{tab:q2-utility}
{\zihao{5}\resizebox{\ifdim\width>\textwidth\textwidth\else\width\fi}{!}{\begin{tabular}{c c c c c c c c c}
\toprule
选点方案 & $\ell$/m & $g$/m & 交会角/度 & 效用/nat & 最差增益/nat & 最坏交会角/度 & $\rho_{90}$/m & $P(\rho_{\min}\le20$ m$)$ \\
\midrule
效用最优 & 840 & $-$543 & 56.4 & 3.302 & 2.640 & 39.5 & 48.4 & 28.3\% \\
最坏交会角最优（对照准则） & 750 & $-$661 & — & — & — & 41.5 & — & — \\
贴近源（理想解） & 1200 & 10 & 89.7 & 1.177 & 0.223 & 0.0 & 600.0 & 4.2\% \\
沿示向线直冲 & 750 & 0 & — & 1.046 & 0.693 & 0.0 & 375.1 & 1.7\% \\
\bottomrule
\end{tabular}}}
\end{table}

\FloatBarrier

\subsection{结果分析与讨论}

本问的提法决定了结论的正确性。若按“已知源位置时最小化定位区域直径”提问，会得到“第二点应贴近源”的结论；但该解对最不利的可能源位置给出接近零的交会角，不能作为保证型策略。在正确提法下，最优第二点位于示向线上的中间位置并向侧向偏移（840 米、$-$543 米），是“压低纵向不确定度”与“保证交会角”两股力量共同作用的结果；其交会角 56.4 度并不接近理想值 90 度，原因在于“保证第二次观测收到信号”把候选点限制在受限区域内，最优解正是两者的折中。

两个准则的结论相互印证，但不互为独立证明（最坏交会角分别为 39.5 度与 41.5 度），说明结论在所测试的准则、先验与网格范围内经验稳健；又因两个准则均为本文选定，不能据此宣称“独立证明”。这一选点相对两类自然选择有实质改善：与贴近源的理想解相比，二次测量后可直接开火的概率由 4.2\% 提高到 28.3\%（约 6.7 倍），$\rho_{90}$ 由 600.0 米降至 48.4 米；即便退到保证可收的可执行解，增益仍达 22.8 倍，可见“保证收到”的代价有限。

“第二次必须收到”并非题面的硬性约束。题面只说“以获得较好的定位效果”，故本文以保证型为主方案，另给出风险型模型（以接收概率与条件期望、或条件风险价值为准则）作为扩展，两者结论不混用。

在线化方面，源的真实距离未知，故在线执行采用“沿示向线逐级复测获取距离代理、再作垂向偏置”的近似做法；复测点不可沿示向线直冲（否则交会角趋于零）；若第二点返回近场信号（距离不超过 5 米）则直接转为光学清除。本章另有两点局限：一是均匀先验只是最大熵基线，不代表源位置的真实分布；二是本章只对照两个准则与 24 个随机场景，且风险型模型仅作扩展，未经实测检验。

\FloatBarrier

\section{问题三模型的建立与求解}

\subsection{思路流程}

问题三要求在源的个数、位置与频道均不可直接读取的条件下，完成搜索、定位与清除。本章把它建模为覆盖约束下的序贯决策问题，分三层处理不同性质的约束：几何层保证“能测到”，负责探测覆盖设计与航点选型；估计层保证“测得准”，在多条带偏示向度上作稳健定位；决策层保证“清得掉、停得对”，决定清除时机与终止条件。三层各自给出独立的保证型结论。


\subsection{数据预处理与可视化}

本章数据分离线构造与在线实测两类。

（1）离线：航点布局与覆盖性质的计算。目标圆域按 20 m 步长离散为 25445 个网格点；对每个网格点计算到最近航点的距离，取全域最大值作为最大盲区；再按“每个航点的 1000 m 圆盘覆盖全域”的条件求解最小集合覆盖，得到覆盖全域所需的最少航点数。

（2）离线：演练器语义的复刻与方案比选。演练测试每局均需人工开启，为此另建离线复现器：由实测日志反演模拟器语义（$r_i$ 为源自身属性、清除仅取决于 20 m 距离、示向度带约 $0.1^\circ$ 量化误差并偶有 $1^\circ$ 量级离群），并按题面规则累加四项耗时。该复现器按四项分解逐局复算全部日志，计时恒等式的相对误差 0.00\%；其平均总时间与冻结版实测相差约 0.5\%（绝对水平口径），可在不消耗演练次数的前提下比较布局与策略；但它假设示向度无噪声，会低估“清除失败”成本，涉及该成本时以实测为准。

（3）在线：模拟器日志的解析与清洗。程序逐条记录每次请求的路径、载荷与响应，据此作四项处理：仅统计被接受的记录（被拒响应的虚拟时间恒为 0）；虚拟时间取最后一次被接受响应中的值，而非本地累加；网络重试保持幂等；逐条区分检测与清除，并注意清除不改变测向机频道，也不产生切换耗时。由相邻两次合法动作的位置差除以 5 m/s 得移动耗时，据此把虚拟时间分解为四项。

（4）可视化。布局与盲区见图~\ref{fig:q3-layout}，虚拟时间构成见附录图 D2。

\begin{figure}[htbp]
\centering
\includegraphics[width=0.54\linewidth]{fig_q3_layout.png}
\caption{问题三单环 15 点建证集的盲区场与巡回里程对比}\label{fig:q3-layout}
\end{figure}

\subsection{模型建立}

（1）几何层：保证可探测集与覆盖设计。记第 $j$ 个航点为 $p_j$，有效接收半径下界为 $r_{\min}=1000$ m。若源位于 $x$，则凡满足

\begin{equation}
	\|p_j-x\|\le r_{\min}
	\label{eq:visible}
\end{equation}

的航点处，无论该源的真实有效接收半径取区间内何值都必定能收到信号。记

\begin{equation}
	P_x=\{p_j:\|p_j-x\|\le r_{\min}\}
	\label{eq:visible-set}
\end{equation}

则覆盖设计要求 $\max_{x\in\bar B(O,1800)}\min_j\|p_j-x\|\le r_{\min}$（保证至少一个航点可测）。

（2）几何层：点数下界——“可检出”不等于“可交会”。仅要求式~\eqref{eq:visible} 成立时，取中心点加半径 1200 m 正六边形的 7 个航点即可覆盖全域，最大盲区约 969 m。但源位于原点附近时，中心点距离为 0 m、周边六点距离恰为 1200 m（超过 1000 m），只有 1 个航点能保证测到，无法交会定位。故本文加严为二重可见要求

\begin{equation}
	\min_{x\in\bar B(O,1800)}|P_x|\ge2
	\label{eq:double-visible}
\end{equation}

（3）几何层：环形布局的解析与选型。由式~\eqref{eq:visible} 与式~\eqref{eq:double-visible}，最简洁的可行布局是单环：$n$ 个航点均匀分布在半径 $R$ 的圆周上（圆心即机器狗初始位置）。盲区约束方面，离所有航点最远的点是圆心，其距离恰为 $R$，故必有 $R\le r_{\min}=1000$ m。边界覆盖约束方面，半径 1800 m 的目标边界上，一点被某个航点覆盖的角半宽为

\begin{equation}
	\Delta\varphi=\arccos\dfrac{1800^2+R^2-r_{\min}^2}{2\cdot1800\cdot R},\qquad n\ge\dfrac{360^\circ}{2\Delta\varphi}
	\label{eq:cover-angle}
\end{equation}

巡回里程方面，沿环依次访问的里程约为

\begin{equation}
	L(R,n)\approx R+(n-1)\cdot2R\sin\dfrac{\pi}{n}\;\xrightarrow[\;n\ \text{较大}\;]{}\;2\pi R
	\label{eq:ring-length}
\end{equation}

式~\eqref{eq:cover-angle} 与式~\eqref{eq:ring-length} 给出结构性结论：巡回里程基本上只由半径 $R$ 决定（约 $2\pi R$），与点数 $n$ 关系很弱；而 $R$ 被“圆心处必须可测”钉在 $R\le r_{\min}$，故最优半径取在约束边界 $R\to1000$ m。据此在离线上扫描 $R\in[900,1000]$ m、$n\in[12,24]$，得最优 $R=990$ m、$n=15$（为圆心处留 10 m 余量）：最大盲区 990.0 m，全域满足式~\eqref{eq:double-visible}，$\Delta\varphi=25.4^\circ$ 远大于 $360^\circ/(2n)=12^\circ$，其巡回里程约 6753 m，仅为原 17 点建证集（13710 m）的一半。

（4）几何层：频道的覆盖证书。由于每个频道至多对应一个源，可对每个尚未清除的频道单独建证。设该频道上已执行检测且全部返回无信号的航点集合为 $P_c$，若

\begin{equation}
	\bigcup_{p\in P_c}\bar B(p,r_{\min})\supseteq\bar B(O,1800)
	\label{eq:certset}
\end{equation}

则该频道在目标域内不存在未清除的全向源。这就是“频道 $\times$ 空间”笛卡尔积证书，频道维（逐频道建证）与空间维（式~\eqref{eq:certset}）缺一不可。

（5）估计层：保证型可行集与点估计的分工。由问题一，$m$ 条示向度给出可行集

\begin{equation}
	\mathcal T=\Big(\bigcap_{k=1}^{m}W(S_k,\theta_k,1^\circ)\Big)\cap\bar B(O,1800)
	\label{eq:feasible-set}
\end{equation}

与第六章的衔接在于：第六章给出的是离线条件下可自由选点的最优第二检测点准则；在线执行时源位置未知、补测点须与扫描路线配合，故实际采用其简化启发式——清除失败后换到法向偏移的观测点复测，再以最小试探点完成闭环；该准则的完全在线化列入第十章改进方向。

由于每个测点的偏差有界，真源必落在每一条示向带内，故 $\mathcal T$ 恒非空，这是最坏情形保证的基础。实际执行中另取一个点估计用于导航：先对同一频道的多条示向度两两求交（自动剔除近乎平行的射线对与落在域外的交点），取交点坐标的分量中位数作初值，再用 Huber 加权最小二乘\upcite{ref:Huber,ref:Soni2026}精化。须注意，当观测恰有两条不平行示向线时，最小二乘目标可在交点处取零，权重被约去，估计值即两线交点，此时决定精度的是布点是否把交会角做对；稳健精化只在观测数达到三次以上时才起作用。

（6）决策层：清除判据。由问题一，清除成功仅取决于开火点到源的距离不超过 20 m，故存在保证成功的开火点当且仅当

\begin{equation}
	\rho(\mathcal T):=\min_z\max_{x\in\mathcal T}\|z-x\|\le20 m
	\label{eq:rho-t}
\end{equation}

又因最远点距离关于 $x$ 为凸函数，其最大值在凸包\upcite{ref:CG,ref:Andrew}上取得，故 $\rho(\mathcal T)=\rho(\operatorname{conv}\mathcal T)$，可先取凸包再求最小外接圆。按此把开火分为三档：$\rho\le20$ m 直冲最小外接圆圆心；$20<\rho\le40$ m 沿可行集主方向单向偏置 20 m；$\rho>40$ m 继续补测。

（7）决策层：清除时机——最优绕道点。在固定巡回下，“一测到就立即动身去清除”会产生可观的绕道：实测 9 局平均每局移动 6829 s，而巡回本身只有 $13710/5=2742$ s，约 60\% 的移动时间消耗在绕道上。为此引入绕道代价判据：设当前航点为 $w_i$、下一个待访航点为 $w_{i+1}$、待清除源的点估计为 $e$，则该段顺路清除的额外代价为

\begin{equation}
	c_i(e)=\|w_i-e\|+\|e-w_{i+1}\|-\|w_i-w_{i+1}\|
	\label{eq:detour}
\end{equation}

遍历尚未访问的各段取 $c(e)=\min_{j\ge i}c_j(e)$；只有当 $c_i(e)\le c(e)$（当前段即为绕道最省的段）时才在第 $i$ 段顺路清除，否则挂起待后续航点处理。此处须对式~\eqref{eq:detour} 作一处实现层修正：候选段必须限定在证书段之内——若把永不执行的备用航点也纳入候选，就会出现“为省一段永远不会走的绕道而拖延到收尾”的判据越界（实测 Q4-6 因此有 10 个源落到收尾段）。

（8）决策层：终止条件。保证型终止有两条：① 已清除数达题面上界 16；② 建证集（单环 15 点）已全部访问，且已检出源全部清除——由式~\eqref{eq:certset}，此时任何仍存在的源都必已在此前的航点上被检出，故不存在漏源。已定位但尚未清除的源在终止前统一执行收尾清除（按最近邻顺序依次前往）。“连续若干步无新发现”单独列为非保证型的加速规则，不与保证型条件并列。

（9）目标与约束。优化目标为虚拟时间

\begin{equation}
	\min T_{\rm virt}=\underbrace{\sum\dfrac{\text{移动距离}}{v}}_{\text{移动}}+\underbrace{5n_{\rm meas}}_{\text{检测}}+\underbrace{n_{\rm sw}}_{\text{切频}}+\underbrace{(5n_{\rm ok}+3n_{\rm miss})}_{\text{清除}},\qquad v=5\ \text{m/s}
	\label{eq:objective}
\end{equation}

其中 $n_{\rm sw}$ 为相邻两次检测频道不同的次数。须注意，虚拟时间 $T_{\rm virt}$ 是优化目标，程序真实运行时间 $T_{\rm real}$（上限 20 分钟）是硬约束，二者推进速度不成比例，须分别记账。

\subsection{模型求解}

求解算法如下：

\begin{lstlisting}[style=codeblock]
输入：单环建证集 W（R=990 m，n=15，共 15 点），频道 1..20，源数下界 10、上界 16
1  enter；清除集 C <- {}，观测集 O <- {}，挂起集 P <- {}
2  for 每个航点 w_i  in  W（按最近邻排序）:
3      对每个未清除、且不在“已定位挂起”状态的频道 ch:
4          测 ch；若返回示向度则记录 (w_i, theta) 并尝试定位（>=2 条）
5          若可定位则写入挂起集 P（本步不立即前往）
6      按绕道代价判据判断：对每个 e  in  P，若第 i 段是该源的绕道最省段，则顺路前往清除
7      若清除失败：沿示向度方向延伸试探，仍失败则移到法向观测点复测并重新定位
8      若 |C| >= 16：终止
9  建证集访问完毕：对剩余 P 执行收尾清除（最近邻次序）
10 exit；报告清除比例与平均定位清除时间 T_clear 均值 = T_clear / n_cleared（均按虚拟时间计）
\end{lstlisting}

（a）布局选型。在离线复现器上以 400～800 局同种子比选，单环 990 m$\times$15 最优：最大盲区 990 m、巡回里程 6753 m，平均总虚拟时间约 4039 s，清除率 100\%。


（b）策略开关对照。在单环 990 m $\times$ 15 下，以 400 局同种子对六种配置作对照，结果见附录表 D9。


（c）停测门槛与清除时机的联合标定。停测与动身清除须共用同一门槛（累计 $\ge$3 条示向度）：两条射线没有冗余，一旦其一为 $1^\circ$ 量级的离群值，或两线近于平行（夹角仅 $1^\circ\sim3^\circ$），交点误差可达数十米。改用累计 $\ge$3 条并配合 Huber 加权最小二乘可抗 1 条离群，但三条本身近共线时仍无法恢复横向约束；据此把清除试探点由 9 个（仅轴向）扩为 13 个（含对角），并在失败后沿示向度方向延伸到 60 m、必要时换到法向观测点复测，问题三最快一局降至 263 s/源。

（d）演练与正式测试实测。最终配置下 16 局全部 100\% 清除（逐局结果见附录表 D2），其中 16 源案例（5 局）的平均耗时 311 s/源；随后在模拟器正式测试的 3 次机会中三局全部 100\% 清除，平均定位清除时间 395.0 s/源，与演练同量级（表~\ref{tab:q3-formal}）。



正式测试过程见图~\ref{fig:q3-path-progress}（左：第 2 局行走轨迹与清除位置；右：三局清除进度与无效尾段）。

\begin{figure}[htbp]
\centering
\begin{minipage}[t]{0.44\linewidth}\centering
\includegraphics[width=\linewidth]{fig_q3_path.png}
\end{minipage}\hfill
\begin{minipage}[t]{0.54\linewidth}\centering
\includegraphics[width=\linewidth]{fig_q3_progress.png}
\end{minipage}
\caption{左：问题三正式测试第 2 局（13 源）行走轨迹与清除位置；右：三次正式测试的清除进度与无效尾段}\label{fig:q3-path-progress}
\end{figure}


（e）极端布局压力测试。把问题三建证集由 15 点削弱到 4 点（每 $90^\circ$ 一点），逼出“只有一个航点能看到某源”的极端情形，过程见附录图 D5：单观测主动补测分支触发 6 次，全部成功定位并清除；其余源在收尾阶段按最近邻顺序清除，最终 11/11=100\%、虚拟时间 7117.0 s。

三次正式测试的结果按题面表 1 的格式报告，见表~\ref{tab:q3-formal}；三局均以覆盖证书终止，清除率 100\%，平均定位清除时间 395.0 s/源。

\begin{table}[ht]
\centering
\caption{问题三的 3 次正式测试结果（按题面表 1 格式）}\label{tab:q3-formal}
{\zihao{5}\resizebox{\ifdim\width>\linewidth\linewidth\else\width\fi}{!}{\begin{tabular}{c c c c}
\toprule
测试案例编码 & 清除干扰源个数 & 平均定位清除时间 /(s/源) & 程序运行时间 /s \\
\midrule
6A7K-XC7T-GK8K-7GGK & 11 / 11 & 435.7 & 1.6 \\
UEVE-DGYC-PJJ2-VFHY & 13 / 13 & 331.9 & 2.5 \\
CGP2-9BMC-G7MB-HR8D & 11 / 11 & 417.4 & 1.7 \\
合计 / 平均 & 35 / 35 & 395.0 & 1.6$\sim$2.5 \\
\bottomrule
\end{tabular}}}
\end{table}

（f）布局的结构性说明。单环 990 m$\times$15 建证集最大盲区 990 m、全域满足二重可见，其覆盖全域所需最少航点为 8 点——但“证书点数最少”与“巡回里程最短”是两个不同目标：该 8 点散布在各方向，若前置为巡回前缀会把行程从 6753 m 拉长到 11836 m，故本文按环序扫描、仅在绕道最省处顺路清除（选址覆盖与巡回里程的这一权衡即经典的最大覆盖选址问题\upcite{ref:MCLP}）。






（17 点建证集坐标见支撑材料。）

\subsection{结果分析与讨论}

（1）“能检出不等于能交会”是本问的几何关键：7 点布局虽已满足“至少一个航点可测”，但在圆心附近只有 1 个航点可测、无法交会，加严为式~\eqref{eq:double-visible} 的二重可见后才真正可用。这也是问题四中 L1/L2/L3 三级证书的前身。

（2）环形布局的结构性结论把约 2 倍的里程差距（13710 m 降至 6753 m）解释为一条解析事实（式~\eqref{eq:cover-angle} 与式~\eqref{eq:ring-length}），而非参数试错的结果。

（3）移动是主导成本，“何时去清除”比“怎么走”的优化空间更大。9 局实测中移动项约占总时间 85\%，而巡回本身的理论值仅 2742 s（逐项见附录表 D6）；引入式~\eqref{eq:detour} 的最优绕道点判据后，移动与总时间分别下降约 45\% 与 34\%。

（4）覆盖证书的作用是把“扫完即停”由经验观察的权宜规则升级为可证明的保证型终止（式~\eqref{eq:certset}）。



\FloatBarrier

\section{问题四模型的建立与求解}

\subsection{思路流程}

问题四在问题三的基础上增加观测可达性约束：定向源仅在朝向两侧各 $90^\circ$ 的扇区内可被检出，故某点测不到它可能仅因该点落在扇区之外\upcite{ref:Kosa2024}；定向源的负观测不含距离信息，问题三依靠负观测构造覆盖证书并提前终止的做法因此不再成立。本章把策略重构为三层：可见性层给出全域（位置 $\times$ 朝向）保证检出证书，替换问题三依赖负观测的终止逻辑；定位层沿用问题三的两点交会与稳健估计内核；执行层与问题三共用清除判据与计时口径。核心增量是可见性层，把提前终止由启发式规则升级为可复核的证书。


\subsection{数据预处理与可视化}

（1）接收半径的保守化。题面规定每个源的有效接收半径 $r_i\in[1000,1500]$ m 且为源自身属性，取值不可读取；若取某具体值作判据，一旦真实值更小就会漏检。故一切覆盖结论统一以下界 $r_{\min}=1000$ m 为口径，实际实现再留 10 m 数值余量（取 990 m）。相应地，正观测只能断言“距离不超过 1500 m”，而定向源的负观测不提供任何距离信息。

（2）角隙场与凸包等价判据的计算。在目标圆盘的 20 m 网格上，对每个网格点 $x$ 筛出满足 $\|p_j-x\|\le r_{\min}$ 的航点集合 $P_x$，将其相对 $x$ 的方位角排序，取最大相邻循环间隙作为最坏角隙 $\Delta_{\max}(x)$；同一网格上并行记录 $|P_x|$（二重可见）与可见点对的最大交会角（几何可定位）。

（3）全域（位置 $\times$ 朝向）覆盖统计与网格口径。由于“能否被看到”同时取决于源的位置与朝向，统计单元取 $\bar B(O,1800)\times S^1$ 上的位置-朝向对。\textbf{全文网格口径统一如下}：证书集求解以 30 m $\times$ $3^\circ$ 粗筛（检出覆盖）并用 20 m 网格选点，复核加密至 10 m $\times$ $1^\circ$；L2/L3 可定位性审计用 20 m $\times$ $10^\circ$；最终以 5 m $\times$ $1^\circ$ 网格完成最严复核（407073 个位置 $\times$ 360 个朝向 = 1.465 亿对）。

（4）可视化。证书的结构（内环约 1000 m、外环约 1900 m）与覆盖率随航点数的变化见图~\ref{fig:q4-cert-structure}。

\begin{figure}[htbp]
\centering
\includegraphics[width=0.54\linewidth]{fig_q4_cert_structure.png}
\caption{问题四证书的结构与“必检出—可定位”两级覆盖}\label{fig:q4-cert-structure}
\end{figure}

（5）航点布局的构造。题面允许提交目标区域外的坐标，故在域外增设环绕带以补足相对源的方位跨度。问题三的布局历经 19 点、17 点、单环 990 m$\times$15 三代（19 点为早期方案，17 点为交会友好建证集，最终采用单环 990 m$\times$15）；其中早期 19 点布局在问题四下存在漏检（见 8.5（1）），本文据此重新求解为 81 点巡回：前 31 点为全域保证检出证书集，其后 50 点为 L2/L3 补给点。

\subsection{模型建立}

（1）定向源的覆盖模型。设定向方向为 $g\in S^1$，则航点 $p$ 对位置 $x$ 的源可见当且仅当

\begin{equation}
	\|p-x\|\le r_{\min}\ \text{且}\ |\operatorname{angle}(p-x)-g|\le90^\circ
	\label{eq:directed-visible}
\end{equation}

式~\eqref{eq:directed-visible} 表明：对定向源，一次“测不到”既可能是距离超出，也可能是落在覆盖角之外，两种原因不可分辨。这是本章与第七章最本质的差别。

（2）L1 一重可见（可检出）。记 $P_x=\{p_j:\|p_j-x\|\le r_{\min}\}$，其方位角按圆周排序为 $\varphi_1<\cdots<\varphi_m$，循环角隙为

\begin{equation}
	\Delta_j=\begin{cases}\varphi_{j+1}-\varphi_j,&j<m\\360^\circ+\varphi_1-\varphi_m,&j=m\end{cases},\qquad\Delta_{\max}(x)=\max_j\Delta_j
	\label{eq:gap}
\end{equation}

命题 8.1　对任意定向方向 $g$，式~\eqref{eq:directed-visible} 至少被一个 $p_j\in P_x$ 满足，当且仅当

\begin{equation}
	\Delta_{\max}(x)\le180^\circ
	\label{eq:l1}
\end{equation}

证明：每个航点在方向 $g$ 上对应一个长度恰为 $180^\circ$ 的闭区间 $I_j=[\varphi_j-90^\circ,\varphi_j+90^\circ]$，式~\eqref{eq:l1} 等价于 $\bigcup_jI_j=S^1$；而 $\bigcup_jI_j\neq S^1$ 当且仅当存在一个角隙大于 $180^\circ$。证毕。

推论 8.2　式~\eqref{eq:l1} 与下列凸包成员判据等价：

\begin{equation}
	\Delta_{\max}(x)\le180^\circ\iff x\in\operatorname{conv}(P_x)
	\label{eq:convexhull}
\end{equation}

因为角隙超过 $180^\circ$ 等价于全部可测航点落在某个开半圆内，由分离定理即得。式~\eqref{eq:convexhull} 把“循环角隙”的计算化简为凸包成员判定，便于在数值上交叉验证。

（3）L2 二重可见（可交会）。两点交会定位至少需要两个可见航点：

\begin{equation}
	|P_x|\ge2,\quad\forall x\in\bar B(O,1800)
	\label{eq:l2}
\end{equation}

（4）L3 几何可定位。还要求可见点对的交会角不病态：

\begin{equation}
	\exists i\neq j:\ p_i,p_j\in P_x,\ |\sin\psi_{ij}(x)|\ge\sin30^\circ
	\label{eq:l3}
\end{equation}

须强调，L3 只是几何充分条件，不排除误差方向不利、频道未切换、路径与时间限制等因素，故统称“几何可定位证书”。

（5）全域保证检出证书与终止判据。式~\eqref{eq:l1} 是对固定 $x$ 的陈述，要用于终止须提升为对巡回前缀的陈述。设巡回应访问的航点序列为 $w_1,w_2,\dots,w_N$，前缀 $P_K=\{w_1,\dots,w_K\}$，若

\begin{equation}
	\forall(x,\theta)\in\bar B(O,1800)\times S^1,\ \exists w\in P_K:\ \|w-x\|\le r_{\min}\ \wedge\ |\operatorname{angle}(w-x)-\theta|\le89^\circ
	\label{eq:certificate}
\end{equation}

则称 $P_K$ 为保证检出证书集，并记 $K^{*}=\min\{K:P_K\ \text{满足式~\eqref{eq:certificate}}\}$。于是有：

命题 8.3（终止的安全性）　若已扫完 $P_{K^{*}}$ 且此前被检出过的源已全部清除，则本次测试中所有存在的干扰源必然已全部被清除。

证明：由式~\eqref{eq:certificate}，任一可能存在的源必在某个 $w\in P_{K^{*}}$ 上被检出并进入检出集合；又“已检出源全部清除”，故不存在未清除的源。证毕。

命题 8.3 中证书建在“位置 $\times$ 朝向”空间，且不能使用负观测，这正是问题四证书集（31 点）显著大于问题三（8 点）的原因。

（6）源数上界的终止加速。题面保证干扰源个数 $n\in[10,16]$，故若已清除源数达到 16 可立即终止；该判据与式~\eqref{eq:certificate} 独立，可直接使用：

\begin{equation}
	n_{\rm cleared}\ge16\ \Longrightarrow\ \text{终止}
	\label{eq:stop16}
\end{equation}

（7）目标函数与约束。优化目标仍为最小化虚拟时间

\begin{equation}
	\min T_{\rm virt}=\sum\dfrac{\text{移动距离}}{v}+5n_{\rm meas}+n_{\rm sw}+(5n_{\rm ok}+3n_{\rm miss}),\qquad v=5\ \text{m/s}
	\label{eq:objective-full}
\end{equation}

其中 $n_{\rm sw}$ 为相邻两次检测频道不同时的切换次数。按式~\eqref{eq:objective-full}，移动项占绝对主导，故长度为 $K^{*}$ 的证书前缀，其里程是证书方案的直接代价。

\subsection{模型求解}

该问题的选点环节是典型的集合覆盖，本文采用贪心策略求解\upcite{ref:Greedy}，其近似保证与次模最大化的 $(1-1/e)$ 界见文献\upcite{ref:Submod}。求解算法如下：

\begin{lstlisting}[style=codeblock]
输入：目标半径 1800，r_min=990，覆盖半角 89°，候选池 R（12 个半径 x 36 方位 ∪ 原布局）
1  在搜索网格 (30 m x 3°) 上对每个 (x,theta) 生成候选航点的覆盖位
2  贪心集合覆盖：每轮取“新增覆盖对数”最大的候选点，直到全域被覆盖
3  反向删除：逐个尝试删除已选点，若仍覆盖则删除，直到不可删
4  复核与补点：在 10 mx1° 网格上复核；若存在未覆盖对，
       按“对未覆盖对增益最大”补点，直到 10 mx1° 网格零漏覆盖
5  最终复核：5 mx1°（407073 位置 x 360 朝向）零漏覆盖
6  排序：对证书集做最近邻 + 2-opt，得到前缀里程最短的访问序
7  输出：证书前缀 + (原布局 - 证书集) 补给点 = 新巡回
\end{lstlisting}

（a）证书求解结果。最终证书集在 5 m $\times$ $1^\circ$ 网格（1.465 亿个位置-朝向对）上零漏覆盖，结构为“内环 13 点（含原点，其余半径 900～1050 m）+ 外环 18 点（半径约 1900 m，其中 4 点为 1500 m，每 $20^\circ$ 一点）”，共 31 点；经 2-opt 排序后证书前缀里程 18884.7 m，达结构下限的 97\%。点数与多尺度复核结果见表~\ref{tab:q4-cert}，31 个证书点坐标见支撑材料。



\begin{table}[htbp]
\centering
\caption{全域保证检出证书集点数与多尺度复核结果}\label{tab:q4-cert}
{\zihao{5}\begin{tabular}{c c}
\toprule
项目 & 值 \\
\midrule
保底接收半径 $r_{\min}$ / m & 990.0 \\
覆盖半角（留余量）/ ° & 89.0 \\
保证检出证书集点数 $K^{*}$ & 31 \\
L2/L3 补充点数 & 50 \\
新巡回总点数 & 81 \\
证书前缀里程 / m & 18884.7 \\
复核 5 m $\times$ 1° 网格未覆盖对数 & 0 \\
\bottomrule
\end{tabular}}
\end{table}

（b）可定位性（L2/L3）全域审计。覆盖证书只保证“必被检出”（L1），不保证“能被交会定位”。本文对全域（位置 $\times$ 朝向）逐个统计证书段中能检出该对的航点数与最佳交会角（20 m $\times$ $10^\circ$ 网格，即 25445 点 $\times$ 36 个朝向），结果见附录表 D5（证书段与全程的 L1/L2 覆盖率见附录表 D4），其解读与由此驱动的单观测主动补测见 8.5（6）。



新巡回的结构与代价见图~\ref{fig:q4-layout}；前缀长度与覆盖能力的关系、以及“每个位置在最不利朝向下所需的最靠后证书序号”（全域最大值恰为 31，即 $K^{*}$ 的几何来源）见附录图 D4。

\begin{figure}[htbp]
\centering
\includegraphics[width=0.54\linewidth,height=0.36\textheight,keepaspectratio]{fig_q4_layout.png}
\caption{问题四航点集——全域保证检出证书前缀与 L2/L3 补给点}\label{fig:q4-layout}
\end{figure}


（c）演练实测见附录表 D3：Q4-10 为反证实验（去掉外环，暴露漏检 4 源），Q4-P1/Q4-P2 为外环削弱到 6 点的压力测试（P1 暴露 VISITED 顺序缺陷、P2 为修复后重跑），其余各局均 100\% 清除（13 源案例 3 局均值 556 s/源）。针对复测中出现的单频道病态交会做执行层加固后，另跑 5 局验证：全部 100\% 清除、0 次清除失败，平均 401～666 s/源（逐局见附录表 D3）。


（d）由时间构成驱动的改进：顺路清除。把每局虚拟时间按式~\eqref{eq:objective-full} 分解可得，“定位后立即清除”的 4 局移动项占 80.0\%～86.7\%，而证书前缀的巡回里程仅 18884.7 m（折合 3777 s），差额几乎全部来自绕道。据此把清除时机改为“在绕道代价最小的航点上顺路清除、证书扫完后再统一收尾”，移动项下降约 36\%，总时间下降约 25\%（时间构成见附录表 D6）；Q4-6 一局曾暴露判据越界瑕疵（把证书之外的备用航点也计入候选），修正后 Q4-7 的 11 个源全部在证书段内清完。

正式测试第 1 局的行走轨迹、证书前缀与清除位置见图~\ref{fig:q4-path}，该局 11 个源全部在证书段内顺路清除，11 次清除调用全部一次命中；第 2 局出现两个示向度不足的弱观测源，其兜底过程见图~\ref{fig:q4-fallback}：频道 12 在仅有一条示向度的情形下主动补测取得第二条示向度后交会定位并清除，频道 14 的两条示向度交会点直接命中源位，在收尾阶段清除成功，两源均在证书扫描完成前完成清除，故三局仍以证书终止。

\begin{figure}[htbp]
\centering
\includegraphics[width=0.54\linewidth,height=0.36\textheight,keepaspectratio]{fig_q4_path.png}
\caption{问题四正式测试第 1 局（11 源）行走轨迹、证书前缀与清除位置}\label{fig:q4-path}
\end{figure}

\begin{figure}[htbp]
\centering
\includegraphics[width=0.86\linewidth]{fig_q4_fallback.png}
\caption{问题四正式测试第 2 局的弱观测源兜底过程：（a）该局全景与两处兜底清除位置；（b）频道 12 由单观测触发主动补测，两条示向度交会点即清除位置；（c）频道 14 的两条示向度交会点命中源位并在收尾阶段清除}\label{fig:q4-fallback}
\end{figure}

三次正式测试的结果按题面表 1 的格式报告，见表~\ref{tab:q4-formal}；三局清除率均为 100\%，平均定位清除时间 568.5 s/源，三局共 49 次清除调用中 37 次命中、37 个源各清除一次，其余 12 次为候选试探落空。

\begin{table}[htbp]
\centering
\caption{问题四三次正式测试结果（按题面表 1 格式）}\label{tab:q4-formal}
{\zihao{5}\resizebox{\ifdim\width>\linewidth\linewidth\else\width\fi}{!}{\begin{tabular}{c c c c}
\toprule
测试案例编码 & 清除干扰源个数 & 平均定位清除时间 /(s/源) & 程序运行时间 /s \\
\midrule
TY66-BXB7-63BJ-M3Z6 & 11 / 11 & 581.1 & 3.1 \\
N3D8-NNPH-3ZCJ-TVQC & 15 / 15 & 504.1 & 3.3 \\
V3VE-CRSQ-5SMB-NHG6 & 11 / 11 & 620.3 & 3.1 \\
合计 / 平均 & 37 / 37 & 568.5 & 3.1$\sim$3.3 \\
\bottomrule
\end{tabular}}}
\end{table}

\subsection{结果分析与讨论}

（1）“看不到”与“定位不了”是两个层次：L1 只保证存在一个可见航点，交会定位还要求至少两个可见点（L2）且交会角不病态（L3）。19 点布局在 L1 即已失败，而 7 点布局正是在 L2 上大面积失败（10540 个网格点），二者失败的层次不同，改进方向也不同。

（2）覆盖结论是“数值证书”而非“数学保证”：在 20 m 网格、不留余量的口径下，许多布局都能“通过”三级判据（例如最坏角隙 150.0$^\circ<180^\circ$）；但一旦改为“位置-朝向对”，加密网格并留余量，就会出现约 0.01\% 的（位置，朝向）对无法保证检出。故覆盖类结论均须写明网格步长、朝向分辨率与余量；本文采用的 31 点证书在上述最严口径下零漏覆盖。

（3）负观测语义的差别决定了终止逻辑必须重建：定向源的负观测不含距离信息，唯一可用的证据是“该源必在某个航点上被看到”，故证书只能建在位置-朝向空间上，规模由 8 点升至 31 点。

（4）证书的收益是终止可证明安全，代价是证书尾段：Q4-3 中全部源清除于 8552.5 s，证书扫完于 10460.5 s，尾段 1908 s 约占总时间的 18\%；鉴于正式测试每问仅有 3 次机会，以约 18\% 的时间换取“绝不漏检”是值得的。

（5）移动仍是唯一主导成本，布局与巡回是最有效的优化杠杆（时间构成见附录表 D6）：证书前缀里程 18884.7 m 已达结构下限的 97\%，说明“把保证条件写清楚”与“执行更快”并不矛盾。

（6）可定位性审计与单观测主动补测的细节见 9.3：证书段 L2、L3 为 96.18\%、92.70\%（20 m$\times$$10^\circ$ 审计网格口径，细网格见附录表 D4），补足补给点后升至 99.89\%、99.37\%，残余约 0.11\% 的（位置，朝向）对只有一个可见点，由 8 个候选偏移的主动补测兜底（首次尝试即把交会角由 $8^\circ$ 提升到 $64.3^\circ$，几何原理与压力测试见附录图 D6）；离线比选中关闭该机制后总时间反升约 2.6\%。

（7）极端情形压力测试与三处实现缺陷的修复见 9.3。


\FloatBarrier

\section{模型的分析与检验}

\subsection{思路流程}

本章按“模型—算法—策略”三层检验，模型层检验判据与拟合模型本身，用解析对照、蒙特卡洛与反例构造；算法层检验求解器数值正确性，用暴力对照、多起点与残差诊断；策略层检验策略的可重复性与稳健性，用多局演练、时间构成分解与参数扰动。


\subsection{检验设置与数据}

（1）验证网格。目标圆盘统一取 20 m 网格（25445 格点），覆盖类结论另在 10 m、5 m 网格复核；问题四的覆盖单元为“位置 $\times$ 朝向”，朝向分辨率粗筛取 3°、复核取 1°；L2/L3 可定位性审计另用 10°。

（2）演练数据。本文统计以最终配置样本为准：问题三 16 局、问题四 18 局，共 34 局（含病态交会加固后的复测 7 局）。逐局结果见附录表 D2、D3；演练案例空间覆盖见附录表 D1 与图 D1。

（3）时间口径校验。把每局虚拟时间独立分解为移动、检测、切频、清除四项：

\begin{equation}
	T_{\rm virt}=\underbrace{\|\Delta\boldsymbol{r}\|/v}_{\text{移动}}+\underbrace{5n_{\rm meas}}_{\text{检测}}+\underbrace{n_{\rm sw}}_{\text{切频}}+\underbrace{(5n_{\rm ok}+3n_{\rm miss})}_{\text{清除}},\qquad v=5\ {\rm m/s}
	\label{eq:time-decomp}
\end{equation}

用离线复现器逐局复算全部日志，式~\eqref{eq:time-decomp} 的累计值与日志 \texttt{virtual\_time\_s} 逐局吻合（相对误差 0.00\%，计时恒等式口径）。

\subsection{模型层与算法层检验}

\textbf{模型层（问题一、二）。} 问题一的覆盖判据与凸包等价（$\rho(\mathcal T)=\rho(\operatorname{conv}\mathcal T)$）由凸性直接得到；旋转卡壳与暴力枚举吻合到 $1.137\times10^{-13}$ m，圆盘求交按 $R=1500$ m、$N=720$ 边形近似的径向误差 0.0143 m。2000 例两点交会构型上，直径中位数 56.83 m、$\rho/(D/2)$ 最大 1.000457，严格判据下不满足覆盖的比例 1.50\%（Wilson 区间见式~\eqref{eq:wilson}），无容差口径升至 11.35\%；反例族可精确达到 Jung 上界 $2/\sqrt3=1.154701$，说明判据非平凡。问题二的双段律拐点为 20.95 m，解析最优第二测点与精确最优解偏差 0.0 m / 0.0°，但其幅度仅具缩放律性质（相关系数 $r=0.821$），故口径为“解析定序、精确定值”。

\textbf{算法层（布局与证书求解）。} 由式~\eqref{eq:double-visible} 可推出盲区最大值出现在圆心且恰等于环半径（故 $R\le r_{\min}$），而 $L(R,n)\approx2\pi R$ 表明巡回里程只由半径决定，最优半径因而取在约束边界：单环 990 m$\times$15 的离线平均总虚拟时间 4039 s、巡回里程 6753 m，并经三局实测确认。问题四的 31 点证书前缀在 5 m$\times$1° 网格（1.465 亿对）上零漏覆盖；覆盖类结论一律附“网格步长 + 朝向分辨率 + 余量”三项口径（详见 8.5（2））。

\textbf{策略层（终止、补测与稳健性）。} 证书只保证必被检出（L1）：证书段 L2、L3 为 96.18\%、92.70\%（20 m$\times$$10^\circ$ 审计网格口径），补足 50 个补给点后升至 99.89\%、99.37\%；残余约 0.11\% 的（位置，朝向）对只有一个可见点，按 13 源估计单局命中概率约 1.4\%，由单观测主动补测兜底（沿“测点$\rightarrow$源”方位取 8 个候选偏移，实测首次尝试即成功，交会角 64.3°）。停测阈值取 2 条时 4202 s 且每局约 5.5 次清除失败，取 3 条为 4114 s，故采用“累计 $\ge$3 条示向度才停测并动身”（附录表 D9）。

\textbf{压力测试与反证实验。} 用去掉外环（此处沿用问题三的 990 m$\times$15 单环作探针）的探针方案实跑一局：对真实的 14 源（5 全向 + 9 定向）只清除了 10 个、漏检 4 个，说明外环（半径大于 1800 m）是保证“任意朝向的源都被检出”所必需的结构。再把问题四外环由 18 点削弱到 6 点、问题三建证集由 15 点削弱到 4 点：单观测补测分支分别触发 2 次（频道 5、频道 18）与 6 次且全部成功，两局均 100\% 清除（附录图 D6、附录图 D5）；该测试同时修复了分支可达性、越界与缺少失败守卫三处实现缺陷。

\subsection{策略层检验：多局演练统计}

（1）清除率。最终配置下共 34 局演练：32 局实现 100\% 清除（源数覆盖 10～16 的整个区间），另 2 局（问题四的 Q4-10、Q4-P1）为有意削弱的反证与压力实验、未达 100\%，不计入清除率；



（2）总时间与源数的关系：总虚拟时间与源数弱负相关（已清除频道被永久跳过），而指标取“总时间 $\div$ 被清除源数”，故源数少时单源均值偏高。

（3）时间构成分解。由式~\eqref{eq:time-decomp} 分解，移动项占绝对主导，见附录表 D6。


由表可见，顺路清除使问题三移动下降 45\%、总时间下降 34\%，问题四分别下降 36\% 与 25\%，代价是检测与切频上升；另一项负结果是跳过已定位频道的重复测量几乎无效果——未出现的频道仍须逐航点测量，这属覆盖证书的内在要求。

（4）由检验驱动的策略改进。附录表 D8 列出各条优化路径的检验结论；本文对每条路径都用离线复现器与实测做同口径对照，只有通过检验者才进入冻结版。


\subsection{灵敏度\upcite{ref:GSA}与误差分析}

（1）口径与接收半径的敏感性。结论随口径变化（20 m 网格无余量时“通过”，改为位置-朝向单元并留余量后不成立）；接收半径取下界 $r_{\min}=1000$ m 是保守口径，凡按 1000 m 判为可见的航点必真实可见，故证书“不遗漏”对 $r_i$ 任意取值成立。

（2）负观测语义与终止代价。全向源的负观测可排除一个 $r_{\min}$ 圆盘，定向源的负观测不含距离信息，这使问题三的“频道 $\times$ 空间”证书（8 点）在问题四失效，须改用位置-朝向证书（31 点）。启发式提前终止可在航点 23 附近结束但有漏检风险；证书式终止须走完 31 点、尾段为 1908 s 与 3389 s（占 18\%～30\%），因每问仅 3 次正式测试机会，以约 18\% 的时间换取“绝不漏检”是值得的。

（3）源数上界与误差量级。题面保证源数 $n\le16$，故清除数达 16 即可立即终止，该判据与覆盖证书是彼此独立的两个充分条件。定位精度主要受两线交会固有误差与示向角量化误差影响：千米距离下 1° 的示向误差约合 17 m 横向偏差，与 20 m 的清除半径同量级，这也是“就近补测”被否决、“顺路清除 + 收尾复测”被采纳的原因。

\subsection{正式测试与演练的一致性}

按题面要求两问各完成 3 次模拟器正式测试，其案例编码、源数与成绩全部取自正式测试运行日志（表~\ref{tab:q3-formal}、表~\ref{tab:q4-formal}）。六局全部 100\% 清除（问题三 35/35、问题四 37/37），六局均以覆盖证书终止；六局共 142 次清除调用，其中 72 次命中且与 72 个已清除源一一对应，即每个源恰好清除一次、无重复清除，另 70 次为候选试探落空。问题三平均 395.0 s/源、问题四平均 568.5 s/源，与最终配置下演练样本的均值 364.6 s/源、555.6 s/源分别相差 8.3\% 与 2.3\%（演练均值按附录表 D2、D3 中给出 s/源 指标的 32 局统计；另有 2 局反证与压力实验的行为指标缺失），说明演练中得到的结论在正式环境中可复现（图~\ref{fig:formal-summary}）。

二者差异主要来自案例抽样而非算法性能：指标取“总虚拟时间 $\div$ 已清除源数”，而总虚拟时间由建证集扫描主导、与源数弱相关，故源数偏少时单源均值偏高——这解释了正式局（11～15 源）均值略高于演练样本的现象。第 2 局还实测了第八章的兜底分支：两个示向度不足的弱观测源分别由主动补测定位与收尾清除完成，该局经 8 次候选试探落空后仍实现 15 个源各清除一次（图~\ref{fig:q4-fallback}），即长尾修复在正式环境中同样生效。

\begin{figure}[htbp]
\centering
\includegraphics[width=0.82\linewidth]{fig_formal_summary.png}
\caption{六次正式测试成绩与最终配置下演练样本的对照（正式为主、演练为辅）}\label{fig:formal-summary}
\end{figure}

\FloatBarrier

\section{模型的评价、改进与推广}

\subsection{模型的优点}

（1）四问被统一在一条方法链上，而非四个彼此独立的模型（整体脉络见第二章末）：四问共用同一套定位内核与清除判据（$\rho\le20$ m），也共用同一套计时口径，因此各问结论可相互校验。关键决策、质量评价与支撑证据的对应关系总览于表~\ref{tab:model-quality}。

\begin{table}[htbp]
\centering
\caption{四问建模的关键决策与建立质量评价}\label{tab:model-quality}
{\zihao{5}\resizebox{\ifdim\width>\textwidth\textwidth\else\width\fi}{!}{\begin{tabular}{c c c c}
\toprule
问题 & 建模的关键决策 & 建立质量评价（客观） & 支撑证据 \\
\midrule
问题一 & 示向带逐行取交，判据取“以 $D$ 为直径的圆能否覆盖定位区域” & 给出了与原问题等价的几何判据，并把“逐点验证”升级为凸包上的等价性结论 & 旋转卡壳 vs 暴力法偏差 $1.137\times10^{-13}$ m；2000 例严格不满足比例 1.50\%（CI [1.05\%,2.13\%]） \\
问题二 & 将第二次测量建模为带接收约束的主动观测点优化 $\max E[\ln(\rho_1/\rho_2)]$ & 把“再测一次选哪里”从经验选点提升为有明确目标函数的优化问题，且给出 minimax 对照准则 & 解析解与精确计算偏差 0.0 m / 0.0°；最优期望效用 3.302 nat \\
问题三 & 由“盲区最大值等于环半径”得 $R\le r_{\min}$，由 $L\approx2\pi R$ 得最优在约束边界 & 布局由解析结构约束确定，而非参数试错；终止条件由覆盖证书给出 & 巡回里程 6753 m；正式 3 局 100\% 清除（35/35，平均 395.0 s/源） \\
问题四 & 建立“位置 $\times$ 朝向”的全域保证检出证书，并分离 L1 检出与 L2/L3 可定位两级 & 把“是否可能漏源”从经验判断升级为可复核的证书，终止条件依赖证书而非启发式 & 31 点证书在 5 m$\times$1° 口径下 1.465 亿对零漏覆盖；去掉外环后实测漏检 4 源；正式 3 局 100\% 清除（37/37） \\
\bottomrule
\end{tabular}}}
\end{table}

（2）终止条件与布局均为“保证型”而非经验规则：终止由覆盖证书给出——扫完证书前缀且已检出源全部清除时必无漏源（去掉外环实测漏检 4 源即其反证，见 9.3）；布局由“盲区最大值等于环半径”与 $L(R,n)\approx2\pi R$ 两条结构结论定在约束边界，单环 $R=990$ m、$n=15$ 使里程由 13710 m 降至 6753 m（$-$51\%），问题四 31 点证书里程 18884.7 m 已达结构下限的 97\%（推导见 7.5、8.3 与 9.3）。

（3）最终效果汇总。问题三、问题四的最终效果汇总于表~\ref{tab:effect}。

\begin{table}[htbp]
\centering
\caption{问题三、问题四的最终效果}\label{tab:effect}
{\zihao{5}\resizebox{\ifdim\width>\linewidth\linewidth\else\width\fi}{!}{\begin{tabular}{c c c}
\toprule
问题 & 指标 & 结果 \\
\midrule
问题三 & 平均定位清除时间 /(s/源) & 395.0（正式三局 331.9$\sim$435.7） \\
问题三 & 清除率与巡回里程 & 正式 3 局 100\%（35/35）；6753 m \\
问题四 & 平均定位清除时间 /(s/源) & 568.5（正式三局 504.1$\sim$620.3） \\
问题四 & 清除率与全域保证检出 & 正式 3 局 100\%（37/37，每个源恰好清除一次）；31 点证书在 1.465 亿位置-朝向对上零漏覆盖 \\
两问 & 终止判据 & 覆盖证书（可复核；去掉外环反证漏检 4 源） \\
两问 & 程序运行时间 / s & 1.6$\sim$3.3（正式六局；窗口上限 1200 s，余量 $\ge$360 倍） \\
\bottomrule
\end{tabular}}}
\end{table}



（4）健壮性经过针对性检验：除常规演练全部 100\% 清除外，压力测试（见 7.4（e）、9.3）把问题三建证集削弱到 4 点、问题四外环削弱到 6 点，逼出“只有一个航点能看到某源”的极端情形，补测分支触发 6 次与 2 次且全部成功。

\subsection{模型的不足}

（1）覆盖证书是数值证书，不是解析证明；策略性能依赖题面上界。证书集求解用 30 m$\times$3° 粗筛、10 m$\times$1° 复核，L2/L3 可定位性审计用 20 m$\times$10°，最终以 5 m$\times$1° 网格复核（1.465 亿个位置-朝向对零漏覆盖），但尚未逐单元给出解析证明；终止判据之一“清除数达 16 即停”使用了题面“源数 $\le$16”的上界，上界改变需重新标定。

（2）示向度误差导致偶发清除失败。示向度带约 0.1° 量化误差并偶有 1° 量级离群，两条射线近于平行时误差被放大到数十米；“累计 $\ge$3 条示向度 + Huber 加权最小二乘”只能抗单条离群，三条本身近共线（夹角十几度）时横向无约束，独立复测中出现过单频道连续落空的长尾，平均 1160.1 s/源。执行层据此扩大试探范围、失败后换法向观测点复测并设落空上限，把该长尾的时间代价压缩约 57\%，残余风险仍属已知项。

（3）离散化与复现器的局限。布局建立在 20 m 网格与 24°/60° 等离散结构上，未在连续域对航点位置作联合最优，而“最小覆盖集”与“最短巡回\upcite{ref:Dijkstra}”是两个不同目标（实测：前置最小覆盖集反使路程由 6753 m 增至 11836 m）。离线复现器中，问题三复现器已校准至 0.5\%（绝对水平口径），问题四复现器把定向方向当作均匀先验，绝对水平比实测保守约 12\%，仅用于相对比较。

\subsection{模型的改进方向}

（1）改进型混合扫描与联合路径规划。当前检测耗时的主体是“未出现的频道须在每个航点上测一遍”（笔数约 $(20-n_{\rm src})\times|W|$），而按覆盖证书只需在最小覆盖子集（问题三 8 点）上测量，即可证明该频道内不存在源；若把该子集前置为巡回前缀并与巡回联合求解（受约束 TSP），检测耗时可再降 30\%～50\%。

（2）把问题二的准则在线化。该准则目前只以简化形式进入执行层（清除失败后沿法向偏移复测）；更优的做法是在可行域半径仍大于清除阈值时按其动态选择下一个测量点，替代固定的候选偏移补测。

（3）覆盖证书的解析化。按细胞分解把目标圆盘划为有限拓扑单元，在每个单元内证明可测航点集合与角隙表达式不变，并用角度传播界 $|\delta\varphi|\le\arcsin\frac{h}{d-h}$ 处理网格误差，即可把“数值证书”升级为解析保证。

（4）自适应接收半径估计。$r_i\in[1000,1500]$ 不可读，本文一律取下界以保证不漏检；若用正、负观测（$d\le r_i$ 与 $d>r_i$）的相容性反推每个源的可行区间，可在不牺牲保证的前提下缩短补测距离。

（5）学习型调度与多机器人协同。GNN / Transformer / Mamba 类方法可用于航点—频道排序，但本文未采用：训练样本仅十余局，且学习型策略给不出覆盖证书，与“绝不漏检”冲突；多机器人方向可把证书扩展为多智能体可见性证书，并实现分区并行搜索与清除，把总时间压到 $\max$ 而非 $\sum$ 量级\upcite{ref:Foster2023}。

\subsection{模型的推广}

（1）问题结构与方法论上的推广。本文处理的是“带朝向约束的可达性 + 巡回清除”类问题：既要求能被观测到，又要求观测后能走到足够近处处置。其中两个动作具有一般性——把“可否被观测”升级为可复核的覆盖证书并使终止依赖它；把“何时处置”改为基于绕道代价的序贯决策（实测带来 25\%～34\% 时间收益）。这种“先建可达性证书、再谈优化”的次序尤其适合漏检代价极高的场景。

（2）具体应用场景。在电子对抗与频谱监测中，无人机对未知定向干扰源（或雷达）的搜索与抵近处置，与问题四同构；灾后应急通信的定向天线对准、植保与消防中继的覆盖巡检路径设计亦对应同一结构。


\section*{AI 工具使用声明}

本参赛队在竞赛过程中使用了 AI 工具（DeepSeek V4.1 Flash），用于代码调试与语言润色，详细情况见支撑材料《AI 工具使用详情.pdf》；全部题面理解、建模思路、模型建立、算法设计、实验方案与结果分析均由参赛队独立完成，AI 参与的少量内容已逐项人工审查核实。

\section*{参考文献}
\renewcommand{\refname}{}
\begin{thebibliography}{12}
\setlength{\itemsep}{0pt}
\setlength{\parsep}{0pt}
\setlength{\topsep}{2pt}
\bibitem{ref:Li2025} Li S J, Lei H T, Zhu C, et al. Bearing-only passive localization and optimized adjustment for UAV formations under electromagnetic silence[J]. Drones, 2025, 9(11): 767.
\bibitem{ref:Torrieri} Torrieri D J. Statistical theory of passive location systems[J]. IEEE Transactions on Aerospace and Electronic Systems, 1984, AES-20(2): 183-198.
\bibitem{ref:Jung} Jung H. Ueber die kleinste Kugel, die eine raeumliche Figur einschliesst[J]. Journal fuer die reine und angewandte Mathematik, 1901, 1901(123): 241-257.
\bibitem{ref:Chen2025} Chen Z, Zhang S J, Fang Y W, et al. The error lower bound of target state estimation in passive localization by bearing-only sensors[J]. IEEE Transactions on Instrumentation and Measurement, 2025, 74: 1-14.
\bibitem{ref:MC} Metropolis N, Ulam S. The Monte Carlo method[J]. Journal of the American Statistical Association, 1949, 44(247): 335-341.
\bibitem{ref:Wilson} Wilson E B. Probable inference, the law of succession, and statistical inference[J]. Journal of the American Statistical Association, 1927, 22(158): 209-212.
\bibitem{ref:Huber} Huber P J. Robust estimation of a location parameter[J]. The Annals of Mathematical Statistics, 1964, 35(1): 73-101.
\bibitem{ref:Soni2026} Soni S, Babu P. Outlier-robust multi-node localization for geoscience applications[J]. Signal Processing, 2026, 248: 110694.
\bibitem{ref:CG} Preparata F P, Shamos M I. Computational Geometry: An Introduction[M]. New York: Springer-Verlag, 1985.
\bibitem{ref:Andrew} Andrew A M. Another efficient algorithm for convex hulls in two dimensions[J]. Information Processing Letters, 1979, 9(5): 216-219.
\bibitem{ref:MCLP} Church R, ReVelle C. The maximal covering location problem[J]. Papers of the Regional Science Association, 1974, 32(1): 101-118.
\bibitem{ref:Kosa2024} K\'osa B, Bukovinszki M, Michaletzky T V, et al. Combined barrier--target coverage for directional sensor network[J]. Sensors, 2024, 24(24): 8093.
\bibitem{ref:Greedy} Chv\'atal V. A greedy heuristic for the set-covering problem[J]. Mathematics of Operations Research, 1979, 4(3): 233-235.
\bibitem{ref:Submod} Nemhauser G L, Wolsey L A, Fisher M L. An analysis of approximations for maximizing submodular set functions---I[J]. Mathematical Programming, 1978, 14(1): 265-294.
\bibitem{ref:GSA} Saltelli A, Ratto M, Andres T, et al. Global Sensitivity Analysis: The Primer[M]. Chichester: John Wiley \& Sons, 2008.
\bibitem{ref:Dijkstra} Dijkstra E W. A note on two problems in connexion with graphs[J]. Numerische Mathematik, 1959, 1(1): 269-271.
\bibitem{ref:Foster2023} Foster A J I, Gianni M, Aly A, et al. Multi-robot coverage path planning for the inspection of offshore wind farms: a review[J]. Drones, 2023, 8(1): 10.
\end{thebibliography}

\newpage
\section*{附录}

\subsection*{附录 A　支撑材料文件列表}

本参赛作品的支撑材料压缩包按 code / figures / tables / latex 四部分组织，共收录 Python 脚本 46 个、航点布局数据文件 1 个、SHA-256 校验清单 1 个（以上合计 48 个文件），另有 LaTeX 源码与附录代码摘录 10 个、插图 24 张（PNG 格式，其中 6 张另附同名 PDF 矢量版）、数据表 13 个；连同本清单与《AI 工具使用详情.pdf》，压缩包共 103 个文件。行为日志由模拟器导出，不计入压缩包。程序文件的结构与用途如表 A1 所示，插图与数据表如表 A2 所示，行为日志如表 A3 所示。

\begin{table}[htbp]
\centering
\caption*{表 A1　支撑材料中的程序与数据文件（46 个脚本 + 1 个布局数据 + 1 个校验清单 + LaTeX 源码）}
{\zihao{5}\resizebox{\textwidth}{!}{\begin{tabular}{c c c c}
\toprule
目录 & 文件数 & 代表文件 & 说明 \\
\midrule
code/01\_冻结执行版 & 2 & run6\_robust.py、waypoints\_p4.json & 主策略（问题三、四通用）与问题四航点布局 \\
code/02\_核心算法代码 & 8 & q1\_q2\_analysis.py、q4\_detect\_cover.py & 问题一、二数值实验，建证集与全域证书求解，L2/L3 审计，离线复现器与时间分解 \\
code/03\_绘图与可视化 & 8 & ch78\_figures.py、ch8\_q4\_figures.py & 论文插图与设计期流程图的绘制脚本 \\
code/04\_其他实验与过程 & 29 & q3\_layout\_opt.py、q4\_prefix\_eval.py & 布局选型、策略开关、消融与被否决方案的验证脚本 \\
latex/ & 10 & main.tex、appendix.tex、appendix\_code\_excerpt/ & 论文与附录的 LaTeX 源码、AI 详情源码及附录代码摘录（编译说明见该目录 说明.md） \\
SHA-256 校验清单 & 1 & \_哈希清单.txt & 冻结执行版两个文件（主策略与航点布局）的 SHA-256 校验值 \\
\bottomrule
\end{tabular}}}
\end{table}

\begin{table}[htbp]
\centering
\caption*{表 A2　支撑材料中的图件与数据表}
{\zihao{5}\resizebox{\textwidth}{!}{\begin{tabular}{c c c}
\toprule
类别 & 数量 & 说明 \\
\midrule
插图（fig\_*） & 24 张 & 与正文图 1$\sim$12 及附录图 D1$\sim$D6 一一对应；均为 PNG 格式，其中 6 张另附同名 PDF 矢量版 \\
数据表（tab\_*） & 13 个 & 正文与附录中三线表与关键插图的源数据，含问题一覆盖统计、问题二效用与鲁棒性、问题三证书与演练汇总、问题四证书与可定位性审计、演练案例空间等 \\
行为日志（.jlog） & 按需附 & 演练按局导出（本文统计样本 34 局，见附录表 D2、D3；案例空间见附录表 D1）与正式测试每问 3 局；正式测试日志不得改名 \\
\bottomrule
\end{tabular}}}
\end{table}

\begin{table}[htbp]
\centering
\caption*{表 A3　行为日志清单（演练按局导出，本文统计样本 34 局；正式测试每问 3 局）}
{\zihao{5}\resizebox{\textwidth}{!}{\begin{tabular}{c c c}
\toprule
类别 & 命名 & 说明 \\
\midrule
正式测试 & formal-p3-*.jlog、formal-p4-*.jlog & 每问 3 局，正式成绩即由这 6 局产生；文件名不得修改 \\
演练测试 & practice-p3-*.jlog、practice-p4-*.jlog & 演练日志按局保存；本文统计样本 34 局（附录表 D2、D3），案例空间覆盖见附录表 D1 \\
\bottomrule
\end{tabular}}}
\end{table}

\subsection*{附录 B　正文中移出与补充的图表}

正文中部分明细表与补充图按“正文只留结论、明细入附录”的方式组织，全部编号为 附录表 D1$\sim$D11 与 附录图 D1$\sim$D6，正文引用处已一一对应。

% ===== 从正文移出的表格（附录用，2026-09-12 重建）=====

\begin{table}[htbp]
\centering
\caption*{表 D1　演练案例空间覆盖（最终配置下的演练辅助样本，含反证与压力实验；正式测试案例见正文表 3、表 5 与附录表 D11）}
{\zihao{5}\begin{tabular}{c c c c c}
\toprule
问题 & 源数 & 全向源数 & 定向源数 & 案例数 \\
\midrule
\midrule
3 & 10 & 10 & 0 & 2 \\
3 & 11 & 11 & 0 & 7 \\
3 & 12 & 12 & 0 & 2 \\
3 & 13 & 13 & 0 & 5 \\
3 & 14 & 14 & 0 & 2 \\
3 & 15 & 15 & 0 & 3 \\
3 & 16 & 16 & 0 & 7 \\
4 & 10 & 3 & 7 & 2 \\
4 & 10 & 9 & 1 & 1 \\
4 & 11 & 0 & 11 & 1 \\
4 & 11 & 10 & 1 & 1 \\
4 & 11 & 2 & 9 & 1 \\
4 & 12 & 11 & 1 & 1 \\
4 & 12 & 7 & 5 & 2 \\
4 & 12 & 8 & 4 & 1 \\
4 & 13 & 11 & 2 & 1 \\
4 & 13 & 2 & 11 & 1 \\
4 & 13 & 8 & 5 & 1 \\
4 & 14 & 10 & 4 & 1 \\
4 & 14 & 5 & 9 & 2 \\
4 & 14 & 6 & 8 & 1 \\
4 & 14 & 8 & 6 & 1 \\
4 & 15 & 11 & 4 & 1 \\
4 & 15 & 14 & 1 & 1 \\
4 & 16 & 5 & 11 & 1 \\
4 & 16 & 8 & 8 & 1 \\
\bottomrule
\end{tabular}}
\end{table}

\begin{table}[htbp]
\centering
\caption*{表 D2　问题三最终配置与对照局的演练记录（辅助样本；单环 990 m$\times$15 建证集 + 顺路清除，含加固后回归复测 2 局）}
{\zihao{5}\resizebox{\ifdim\width>\linewidth\linewidth\else\width\fi}{!}{\begin{tabular}{c c c c c c}
\toprule
局 & 源数 & 清除 & 虚拟时间 / s & 平均耗时 /(s/源) & 终止方式 \\
\midrule
Q3-12 & 10 & 10 & 3981.0 & 398 & 单环 990 m$\times$15 建证集扫完 \\
Q3-13 & 15 & 15 & 4324.5 & 288 & 单环 990 m$\times$15 建证集扫完 \\
Q3-14 & 15 & 15 & 4372.8 & 291 & 单环 990 m$\times$15 建证集扫完 \\
Q3-15 & 11 & 11 & 4306.1 & 391 & 按条数停测=2 \\
Q3-16 & 16 & 16 & 6350.9 & 397 & 源数上界终止 \\
Q3-17 & 15 & 15 & 5189.8 & 346 & 单环 \\
Q3-18 & 16 & 16 & 4207.7 & 263 & 源数上界终止 \\
Q3-19 & 16 & 16 & 4870.1 & 304 & 源数上界终止 \\
Q3-20 & 12 & 12 & 4431.2 & 369 & 单环 \\
Q3-21 & 16 & 16 & 4762.1 & 298 & 源数上界终止 \\
Q3-22 & 10 & 10 & 5059.5 & 506 & 建证集扫完 \\
Q3-23 & 13 & 13 & 4485.2 & 345 & 建证集扫完 \\
Q3-P1 & 10 & 10 & 3818.5 & 381.8 & 压力测试：建证集削到 6 点（补测未触发） \\
Q3-P2 & 11 & 11 & 7117.0 & 647 & 压力测试：建证集削到 4 点（补测触发 6 次均成功） \\
P3-R1 & 16 & 16 & 4705.8 & 294 & 加固后回归复测（延伸试探/法向复测/空清上限） \\
P3-R2 & 14 & 14 & 4389.3 & 314 & 加固后回归复测 \\
\bottomrule
\end{tabular}}}
\end{table}

\begin{table}[htbp]
\centering
\caption*{表 D3　问题四最终配置与对照局的演练记录（辅助样本；31 点证书 + 顺路清除 + 单观测补测，含反证与压力实验、加固后复测 5 局）}
{\zihao{5}\resizebox{\ifdim\width>\linewidth\linewidth\else\width\fi}{!}{\begin{tabular}{c c c c c c}
\toprule
局 & 源数 & 清除率 & 虚拟时间 / s & 平均耗时 /(s/源) & 策略版本 \\
\midrule
Q4-5 & 15 & 100\% & 8339.5 & 556.0 & 31 点证书 + 顺路清除 \\
Q4-6 & 14 & 100\% & 8560.3 & 611.4 & 31 点证书 + 顺路清除 \\
Q4-7 & 11 & 100\% & 7610.6 & 691.9 & 31 点证书 + 顺路清除（绕道判据修正） \\
Q4-8 & 13 & 100\% & 7061.0 & 543.2 & 累计 $\ge$3 条示向度才开清（零复测） \\
Q4-10 & 10 / 14 & 71\% & 5397.9 & — & 反证实验：去掉外环（仅内环 990 m$\times$15），漏检 4 源 \\
Q4-11 & 15（11 全向 + 4 定向） & 100\% & 6254.2 & 416.9 & 全向多于定向的构成下最快一局 \\
Q4-12 & 16 & 100\% & 6541.3 & 408.8 & 加入单观测主动补测 \\
Q4-13 & 12 & 100\% & 8575.0 & 714.6 & 补测阈值放宽到 $\le$2（正常流程未触发） \\
Q4-14 & 11（10 全向 + 1 定向） & 100\% & 6542.0 & 594.7 & 补测候选扩到 8 个 \\
Q4-15 & 13 & 100\% & 7756.6 & 596.7 & 冻结版 \\
Q4-16 & 14 & 100\% & 6628.8 & 473.5 & 冻结版冒烟 \\
Q4-P1 & 8 / 11 & 73\% & 5976.9 & — & 压力测试：外环人为削弱到 6 点（暴露 VISITED 顺序缺陷） \\
Q4-P2 & 12 & 100\% & 8697.8 & 724.8 & 压力测试：修复后重跑（补测触发 2 次且均成功，见 9.3） \\
P4-R1 & 15 & 100\% & 6563.7 & 437.6 & 加固后复测（延伸试探 + 法向复测） \\
P4-R2 & 13 & 100\% & 6877.7 & 529.1 & 加固后复测 \\
P4-R3 & 15 & 100\% & 7854.8 & 523.7 & 加固后复测 \\
P4-R4 & 12 & 100\% & 7987.6 & 665.6 & 加固后复测（含单频道空清上限） \\
P4-R5 & 16 & 100\% & 6420.6 & 401.3 & 加固后复测（含单频道空清上限） \\
\bottomrule
\end{tabular}}}
\end{table}


\begin{table}[htbp]
\centering
\caption*{表 D4　证书前缀与全程的可定位性（5 m$\times$1$^\circ$ 口径）}
{\zihao{5}\begin{tabular}{c c c c}
\toprule
巡回前缀 & L1 保证检出率 & L2 两点可定位率 & 全域最坏角隙 \\
\midrule
新 31 点证书前缀 & 100.00\% & 95.56\% & — \\
新 81 点全程 & 100.00\% & 99.82\% & 127.7° \\
\bottomrule
\end{tabular}}
\end{table}

\begin{table}[htbp]
\centering
\caption*{表 D5　可定位性（L2/L3）全域审计（20 m$\times$10$^\circ$ 网格口径；细网格 5 m$\times$1$^\circ$ 的 L1/L2 见表 D4）}
{\zihao{5}\resizebox{\ifdim\width>\linewidth\linewidth\else\width\fi}{!}{\begin{tabular}{c c c c c}
\toprule
航点集合 & 累计航点数 & L1 必检出 / \% & L2 可见点$\ge$2 / \% & L3 交会角$\ge$30° / \% \\
\midrule
内环 13 点 & 13 & 78.05 & 64.66 & 46.03 \\
证书段（内环 13 + 外环 18） & 31 & 100.0 & 96.18 & 92.7 \\
证书段 + 前 19 个补给点 & 50 & 100.0 & 97.93 & 95.69 \\
证书段 + 全部 50 个补给点 & 81 & 100.0 & 99.89 & 99.37 \\
\bottomrule
\end{tabular}}}
\end{table}




\begin{table}[htbp]
\centering
\caption*{表 D6　问题三与问题四对照组的虚拟时间构成分解（按式~\eqref{eq:time-decomp} 分解）}
{\zihao{5}\resizebox{\ifdim\width>\linewidth\linewidth\else\width\fi}{!}{\begin{tabular}{c c c c c c c}
\toprule
策略版本 & 虚拟时间 / s & 移动 / s & 检测 / s & 切频 / s & 清除 / s & 移动占比 \\
\midrule
问题三：立即跑去清除（9 局平均） & 8046 & 6829 & 941 & 187 & 89 & 84.9\% \\
问题三：顺路清除（Q3-9～Q3-11 平均） & 5300 & 3765 & 1212 & 241 & 83 & 71.0\% \\
问题三：单环 990 m$\times$15（Q3-12、Q3-13 平均） & 4153 & 2881 & 983 & 196 & 94 & 69.4\% \\
问题四：立即跑去清除（Q4-3、Q4-4 平均） & 10962 & 8931 & 1568 & 310 & 154 & 81.5\% \\
问题四：顺路清除（Q4-5～Q4-7 平均） & 8170 & 5691 & 2010 & 401 & 69 & 69.7\% \\
\bottomrule
\end{tabular}}}
\end{table}

\begin{table}[htbp]
\centering
\caption*{表 D7　问题二基准场景的几何量与补充数值}
{\zihao{5}\resizebox{\ifdim\width>\linewidth\linewidth\else\width\fi}{!}{\begin{tabular}{c c c}
\toprule
量 & 数值 & 说明 \\
\midrule
基准$d_1$ / m & 1200 & $S_1=(0,0)$、$G=(0,1200)$ \\
单点观测$D_1$ / m & 1500.06 & 仅含径向不确定度 \\
最优$S_2$ & (10.5, 1200.0) & — \\
最优$d_2$ / m & 10 & — \\
最优交会角 / 度 & 89.7 & — \\
$D_{\min}$ / m & 31.44 & 精确几何 \\
相对单次增益 & 47.7 倍 & $D_1/D_{\min}$ \\
解析与数值最优点偏差 & 0.0 m / 0.0 度 & — \\
垂线偏移拐点$a^{*}$ / m & 20.95 & $a^{*}=d_1\tan1^\circ$ \\
可执行最优$D$ / m & 65.76 & 源位置未知、要求第二点保证可收 \\
可执行最优增益 & 22.8 倍 & $D_1/D$ \\
解析平台值 / m & 41.89 & $2d_1\tan1^\circ$，非最优值 \\
解析 vs 精确相关系数（全网格） & 0.82089 & 含近退化几何 \\
解析 vs 精确相关系数（$d_2\ge200$ m） & 0.82112 & 缩放律适用区 \\
$a=5.5$ m 时 $D$ / m & 26.46 & 受 5 米盲区约束的下界 \\
$a=a^{*}$ 时 $D$ / m & 41.47 & 拐点 \\
$a=200$ m 时 $D$ / m & 42.48 & 平台段 \\
稳健推荐偏移 / m & 41.9 & 落在平台起点，对估计误差不敏感 \\
$R_{\rm recv}$ 定义 & $\Omega_1$ 全部极点对应各 1000 米圆盘之交 & 精确保证可收集；一般不是圆盘 \\
圆盘外近似圆心 / 半径 / m & (0, 750) / 661.44 & 仅作对照，用于保证型优化不安全 \\
可执行最优$\ell$ / m & 920 & 沿示向线前进量 \\
可执行最优$g$ / m & $-$392 & 侧向偏移 \\
可执行最优交会角 / 度 & 54.5 & — \\
\bottomrule
\end{tabular}}}
\end{table}
\begin{table}[htbp]
\centering
\caption*{表 D8　优化路径的检验结论}
{\zihao{5}\resizebox{\ifdim\width>\linewidth\linewidth\else\width\fi}{!}{\begin{tabular}{c c c}
\toprule
优化 & 结论 & 判据 \\
\midrule
建证集替代经验布局（Q3 17 点、Q4 31 点） & 采纳 & 覆盖能力提升且里程下降 \\
证书式终止替代启发式终止 & 采纳 & 终止可证明安全（命题 8.3） \\
顺路清除替代立即清除 & 采纳 & Q3 移动 6829 降至 3765 s、总时间 8046 降至 5300 s；Q4 移动 8931 降至 5691 s、总时间 10962 降至 8170 s \\
单环 990 m$\times$15 替代 17 点交会友好建证集 & 采纳 & Q3 巡回 13710 降至 6753 m、平均 5300 降至 4153 s（离线 5275 降至 4039 s，-23\%） \\
清除试探点由 $\pm$20 m 收紧至 $\pm$18 m & 采纳（对冲） & 实测 $\pm$20 m 试探点偶有失败，疑判据为“严格 <20 m”；$\pm$18 m 对两种判据都安全 \\
\bottomrule
\end{tabular}}}
\end{table}

\begin{table}[htbp]
\centering
\caption*{表 D9　策略开关对照（单环 990 m$\times$15，400 局同种子）}
{\zihao{5}\resizebox{\ifdim\width>\linewidth\linewidth\else\width\fi}{!}{\begin{tabular}{c c c c c}
\toprule
配置 & 平均总时间/s & 移动/s & 检测/s & measure 次数 \\
\midrule
顺路清除 + 累计 $\ge$3 条示向度才停测、才动身清除（本文采用） & 4109 & 2893 & 975 & 192 \\
仅按条数：取得 2 条即停测 & 4202（σ=0.05° 时 5.5 次清除失败/局） & 2893 & 902 & 180 \\
仅按条数：取得 3 条即停测 & 4114 & 2893 & 975 & 193 \\
顺序自适应（纯最近任务，无近旁优惠） & 4566 & 3464 & 865 & 173 \\
覆盖子集前置 + 动态终止 & 5062 & 3826 & 977 & 195 \\
关闭顺路清除（改为立即前往） & 5721 & 3857 & 1500 & 300 \\
\bottomrule
\end{tabular}}}
\end{table}


\begin{table}[htbp]
\centering
\caption*{表 D10　两点交会构型的直径与覆盖判据统计（含极值搜索比值上界）}
{\zihao{5}\resizebox{\ifdim\width>\linewidth\linewidth\else\width\fi}{!}{\begin{tabular}{c c c}
\toprule
指标 & 数值 & 说明 \\
\midrule
采样例数 & 2000 & 两点交会，源位在半径 1700 米内均匀 \\
直径中位数 / m & 56.83 & 定位区域直径 \\
$\rho_{\min}/(D/2)$ 最大值 & 1.000457 & 实测 \\
不满足覆盖比例 & 1.50\% & 阈值$1+10^{-9}$ \\
不满足覆盖比例（含浮点平局） & 11.35\% & 阈值 1，口径敏感性对照 \\
Wilson 95\% 置信区间 & [1.05\%, 2.13\%] & — \\
$\rho_{\min}/(D/2)$ 中位数 & 1.000000 & — \\
极值搜索（2 点 20 万例；3～6 点各 6 万例） & 1.0006～1.0043 & 随机极值上界 \\
最坏构型（三点偏移） & 1.154701 & 反例族 \\
Jung 上界$2/\sqrt3$ & 1.154701 & — \\
旋转卡壳 vs 暴力枚举最大偏差 / m & $1.137\times10^{-13}$ & 算法一致性 \\
\bottomrule
\end{tabular}}}
\end{table}

\begin{figure}[htbp]
\centering
\includegraphics[width=0.62\linewidth]{fig_case_space.png}
\caption*{图 D1　演练案例空间覆盖（源数与全向/定向构成）}
\end{figure}

\begin{figure}[htbp]
\centering
\includegraphics[width=0.62\linewidth]{fig_q3_time.png}
\caption*{图 D2　六次正式测试的虚拟时间构成分解（全部取自正式测试运行日志）}
\end{figure}

\begin{figure}[htbp]
\centering
\includegraphics[width=0.62\linewidth]{fig_q1_ratio_hist.png}
\caption*{图 D3　覆盖判据的口径敏感性与尾部互补分布}
\end{figure}

\begin{figure}[htbp]
\centering
\includegraphics[width=0.62\linewidth]{fig_q4_cover.png}
\caption*{图 D4　全域保证检出证书的覆盖能力曲线与最坏情形}
\end{figure}

\begin{table}[htbp]
\centering
\caption*{表 D11　六次正式测试运行明细（正式成绩见表 3 与表 5）}
{\zihao{5}\resizebox{\ifdim\width>\linewidth\linewidth\else\width\fi}{!}{\begin{tabular}{c c c c c c c c}
\toprule
局次 & 测试案例编码 & 源数 & 清除 & 虚拟时间 / s & 平均 /(s/源) & 程序运行 / s & 请求数 \\
\midrule
\midrule
问题三·第 1 局 & 6A7K-XC7T-GK8K-7GGK & 11 & 11 & 4792.8 & 435.7 & 1.6 & 256 \\
\midrule
问题三·第 2 局 & UEVE-DGYC-PJJ2-VFHY & 13 & 13 & 4314.4 & 331.9 & 2.5 & 225 \\
\midrule
问题三·第 3 局 & CGP2-9BMC-G7MB-HR8D & 11 & 11 & 4591.6 & 417.4 & 1.7 & 253 \\
\midrule
问题四·第 1 局 & TY66-BXB7-63BJ-M3Z6 & 11 & 11 & 6392.4 & 581.1 & 3.1 & 368 \\
\midrule
问题四·第 2 局 & N3D8-NNPH-3ZCJ-TVQC & 15 & 15 & 7561.7 & 504.1 & 3.3 & 336 \\
\midrule
问题四·第 3 局 & V3VE-CRSQ-5SMB-NHG6 & 11 & 11 & 6823.4 & 620.3 & 3.1 & 351 \\
\bottomrule
\end{tabular}}}
\end{table}

六局均以覆盖证书终止（问题三为建证集 15 点扫完，问题四为证书前缀 31 点扫完），
六局共 142 次 \texttt{/clear} 调用，其中 72 次命中、与 72 个已清除源一一对应（每个源恰好清除一次，无重复清除），另 70 次为候选试探落空（\texttt{no\_target\_in\_range}）；第 2 局的兜底分支触发情况见正文图~\ref{fig:q4-fallback}。

\begin{figure}[htbp]
\centering
\begin{minipage}[t]{0.46\linewidth}\centering
\includegraphics[width=\linewidth]{fig_q3_stress.png}
\end{minipage}\hfill
\begin{minipage}[t]{0.46\linewidth}\centering
\includegraphics[width=\linewidth]{fig_q3_stress_progress.png}
\end{minipage}
\caption*{图 D5　问题三压力测试：左为稀疏建证集逼出的单观测情形，右为建证集走完后由补测与收尾完成清除}
\end{figure}

\begin{figure}[htbp]
\centering
\begin{minipage}[t]{0.48\linewidth}\centering
\includegraphics[width=\linewidth]{fig_single_fix.png}
\end{minipage}\hfill
\begin{minipage}[t]{0.48\linewidth}\centering
\includegraphics[width=\linewidth]{fig_stress.png}
\end{minipage}
\caption*{图 D6　左：单观测主动补测的几何原理；右：问题四压力测试——削弱外环逼出单观测情形}
\end{figure}

\FloatBarrier

\subsection*{附录 C　核心源程序代码}

附录按“正文只留结论、代码只留核心”的方式收录主策略程序与各问核心算法的关键函数，完整可运行源程序（46 个脚本）随支撑材料提供。复现环境为 Python 3.12，仅使用标准库与 NumPy、SciPy、Matplotlib；策略进程与模拟器通过本地 HTTP 接口交互，无外部依赖、无训练过程。运行方式：设置环境变量 \texttt{CUMCM\_ROBOT\_ID} 后执行 \texttt{python run6\_robust.py p3}（问题四改为 p4）。代码中的本机路径已替换为占位说明，参赛队号以环境变量注入。附录代码均为参赛队自研实现，未复制任何第三方源码；最小外接圆、凸包与旋转卡壳、Huber 加权最小二乘等经典算法的通用参考见参考文献\upcite{ref:Jung,ref:Huber,ref:CG,ref:Andrew}。

\subsubsection*{C.1　主策略程序 run6\_robust.py（核心函数摘录）}
\lstinputlisting[style=codeappx]{appendix_code_excerpt/run6_robust.py}

\subsubsection*{C.2　问题四全域保证检出证书 q4\_detect\_cover.py（核心函数摘录）}
\lstinputlisting[style=codeappx]{appendix_code_excerpt/q4_detect_cover.py}

\subsubsection*{C.3　问题三建证集 q3\_certset.py（核心函数摘录）}
\lstinputlisting[style=codeappx]{appendix_code_excerpt/q3_certset.py}

\subsubsection*{C.4　问题四 L2/L3 可定位性审计 q4\_l2l3\_audit.py（核心函数摘录）}
\lstinputlisting[style=codeappx]{appendix_code_excerpt/q4_l2l3_audit.py}

\subsubsection*{C.5　问题一、二数值实验 q1\_q2\_analysis.py（核心函数摘录）}
\lstinputlisting[style=codeappx]{appendix_code_excerpt/q1_q2_analysis.py}



\end{document}
